
%% ── 4-Capas-Codigo/4-3-Visualizacion.tex
\subsection*{4 3 Visualizacion}
% !TeX root = ../tesis.tex
%%%%%%%%%%%%%%%%%%%%%%%%%%%%%%%%%%%%%%%%%%%%%%%%%%%%%%%%%%%%%%%%%%%
% SECCIÓN 4.3 — Transport-gis: Motor de Visualización
%%%%%%%%%%%%%%%%%%%%%%%%%%%%%%%%%%%%%%%%%%%%%%%%%%%%%%%%%%%%%%%%%%%

% =============================================================================
%  4.3 — VISIÓN GENERAL
% =============================================================================
\section{Transport-gis: El motor de Visualización}
\label{sec:transport-gis-intro}

\textbf{Transport-gis-zmvm-mjg} (\termino{Transport GIS — Zona Metropolitana del Valle de México, M.J. Garibelt}) 
es el tercer y último componente de la arquitectura de código de esta investigación. Mientras \textbf{Apimetro} 
opera como el motor de datos unificado —responsable de la ingesta, estandarización y exposición de la red de transporte 
como servicio— y \textbf{VFTModel} constituye el núcleo analítico que calcula los indicadores de desempeño de la red, 
\textbf{Transport-gis} representa la capa de visualización: el sistema que traduce los resultados del modelo en 
representaciones cartográficas interpretables para el análisis urbanístico.

Al igual que los componentes precedentes, la denominación del sistema proviene de la obra literaria ficcional 
\textit{Saga Diarios, tomo~2}, de autoría del investigador. En dicha novela, la urbanista ficcional 
\textit{Mildred Julietta Cordero Garibelt} —referida por su acrónimo \textit{M.J.G}— es la responsable de la cartografía 
y la planeación espacial del sistema de movilidad de \textit{Ciudad Capital}, ucronía de la Ciudad de México en la década de 1960. 
La trama sitúa a este personaje en el centro de eventos políticos relevantes de la historia mexicana del siglo~\textsc{xx} 
\parencite{lopezmateos2023comoelcantodelcisne}, cuya intervención en las decisiones sobre infraestructura y movilidad urbana articula 
el argumento central de la novela. El paralelismo con la presente investigación es deliberado: así como \textit{M.J.G} produce 
los mapas que hacen legible el sistema de movilidad de \textit{Ciudad Capital}, \textbf{Transport-gis} genera las representaciones 
cartográficas que permiten interpretar los resultados del \textbf{VFTModel} sobre la red metropolitana real de la \zmvm.

\subsection{Transport-gis: Su arquitectura y funcionamiento}
\label{subsec:transport-gis-arquitectura-funcionamiento}

\textbf{Transport-gis} se organiza en cuatro módulos funcionales interdependientes: un módulo de \textbf{análisis} (\texttt{analysis/}), que gestiona la consulta a los servicios 
externos y el procesamiento de indicadores; un módulo de \textbf{datos} (\texttt{data/processed/}), que almacena las capas geoespaciales en formato \termino{GeoPackage} (\texttt{.gpkg}) 
bajo control de Git~\textsc{lfs}; un módulo \textbf{cartográfico} (\texttt{maps/}), que agrupa los proyectos QGIS, las plantillas tipográficas y los archivos de exportación; y un 
módulo de \textbf{orquestación} (\texttt{pipeline/}), que coordina la ejecución secuencial del proceso completo. La Figura~\ref{fig:tg-pipeline} ilustra el flujo de datos 
entre estos módulos y los componentes externos de la arquitectura de la investigación.

% =============================================================================
\subsubsection{Flujo de producción de datos}
\label{subsubsec:tg-flujo-datos}
% =============================================================================

El proceso de producción cartográfica parte de dos fuentes primarias externas: los \textit{feeds} \gtfs{} de la \textsc{semovi} (2024), que contienen los horarios, 
rutas y paradas de los sistemas de transporte de la \cdmx{}; y la cartografía base del \textsc{inegi} (2023), que provee los límites político-administrativos y la 
infraestructura vial de la \zmvm{}. Estas fuentes alimentan a \textbf{Apimetro} —que las transforma en una \api{} \textsc{rest} consumible con la topología del 
grafo de la red— y a \textbf{VFTModel}, que las procesa para calcular los indicadores de eficiencia nodal.

Dentro del repositorio de \textbf{Transport-gis}, el script \texttt{apimetro\_client.py} consulta la \api{} de \textbf{Apimetro} para recuperar la topología del 
grafo y los atributos de cada modo de transporte, mientras que \texttt{vft\_runner.py} ejecuta las rutinas del \textbf{VFTModel} para producir los indicadores por nodo en 
formato \texttt{.csv}. El orquestador \texttt{pipeline/run\_pipeline.py} gestiona la secuencia completa de transformaciones hasta exportar las capas procesadas al 
directorio \texttt{data/processed/}. Los fragmentos de código de estos componentes se \DIFdelbegin \DIFdel{documentan en el Anexo~\ref{anx:transportgis-tecnico}}\DIFdelend \DIFaddbegin \DIFadd{detallarán en un anexo técnico}\DIFaddend .\pendiente{Crear Anexo — Documentación Técnica Transport-gis}

\begin{figure}[htbp]
\centering
\begin{tikzpicture}[
  srcbox/.style = {rectangle, rounded corners=3pt, draw=gray!55, fill=gray!8,
                   minimum width=3.0cm, minimum height=0.90cm,
                   align=center, font=\small},
  extbox/.style = {rectangle, rounded corners=3pt, draw=unamAzul, fill=unamAzul!10,
                   minimum width=3.2cm, minimum height=0.90cm,
                   align=center, font=\small},
  intbox/.style = {rectangle, rounded corners=3pt, draw=unamAzul!60, fill=unamGrisClaro,
                   minimum width=6.4cm, minimum height=1.05cm,
                   align=center, font=\small},
  outbox/.style = {rectangle, rounded corners=3pt, draw=unamOro, fill=unamOro!15,
                   minimum width=3.2cm, minimum height=0.90cm,
                   align=center, font=\small},
  ->, >=Stealth, thick
]
%% — Fuentes externas —
\node[srcbox] (gtfs)  at (0.0, 9.2) {\gtfs{} SEMOVI~2024};
\node[srcbox] (inegi) at (5.4, 9.2) {INEGI~2023 (cartografía)};

%% — Componentes externos —
\node[extbox] (api) at (0.0, 7.5) {\textbf{Apimetro}\\{\footnotesize API REST (Go)}};
\node[extbox] (vft) at (5.4, 7.5) {\textbf{VFTModel}\\{\footnotesize (Python / NetworkX)}};

%% — Transport-gis: análisis —
\node[intbox] (scripts) at (2.7, 5.5)
  {{\footnotesize\texttt{apimetro\_client.py} $\cdot$ \texttt{vft\_runner.py}}\\
   {\footnotesize\texttt{pipeline/run\_pipeline.py}}};

%% — Transport-gis: datos —
\node[intbox] (gpkg) at (2.7, 4.0)
  {{\small\texttt{data/processed/}}\\
   {\footnotesize 10 capas \texttt{.gpkg} (Git~LFS)}};

%% — Transport-gis: QGIS —
\node[intbox] (qgis) at (2.7, 2.5)
  {{\small QGIS~3.x — \texttt{maps/projects/*.qgz}}\\
   {\footnotesize\texttt{maps/templates/*.qpt} (paletas UNAM)}};

%% — Transport-gis: exports —
\node[intbox] (exports) at (2.7, 1.0)
  {{\small\texttt{maps/exports/}}\\
   {\footnotesize\texttt{*.pdf} $\cdot$ \texttt{*.png}}};

%% — Destino final —
\node[outbox] (tesis) at (2.7, -0.5)
  {\textbf{Tesis LaTeX}\\{\footnotesize\texttt{Figures/}}};

%% — Flechas: fuentes → externos —
\draw[color=gray!60] (gtfs)  -- (api);
\draw[color=gray!60] (inegi) -- (vft);

%% — Flechas: externos → scripts (convergentes, con codo exterior) —
\draw[color=unamAzul!70] (api.south) -- ++(0,-0.55) -| (scripts.north west);
\draw[color=unamAzul!70] (vft.south) -- ++(0,-0.55) -| (scripts.north east);

%% — Flujo interno —
\draw[color=unamAzul!70] (scripts) -- (gpkg);
\draw[color=unamAzul!70] (gpkg)    -- (qgis);
\draw[color=unamAzul!70] (qgis)    -- (exports);

%% — Exports → LaTeX —
\draw[color=unamOro] (exports) -- (tesis);

%% — Recuadro del repositorio Transport-gis (en fondo) —
\begin{scope}[on background layer]
  \node[draw=unamAzul!40, dashed, rounded corners=6pt, fill=unamAzul!3,
        fit={(scripts)(gpkg)(qgis)(exports)},
        inner sep=14pt] (tgbox) {};
\end{scope}
%% — Label DENTRO del recuadro, esquina sup-izq —
\node[font=\small\bfseries\color{unamAzul}, anchor=north west]
  at ($(tgbox.north west) + (4pt,-3pt)$) {Transport-gis-zmvm-mjg};

\end{tikzpicture}
\caption[Pipeline de producción cartográfica de Transport-gis]%
  {Flujo de datos del sistema \textbf{Transport-gis-zmvm-mjg}. 
  Los componentes externos \textbf{Apimetro} y \textbf{VFTModel} proveen la información primaria; 
  dentro del repositorio, los \textit{scripts} de análisis consolidan los datos en capas \texttt{.gpkg} 
  que QGIS transforma en los mapas de presentación exportados a la tesis.}
\label{fig:tg-pipeline}
\end{figure}

% =============================================================================
\subsubsection{Capas del sistema de información geográfica}
\label{subsubsec:tg-capas}
% =============================================================================

El módulo de datos mantiene diez capas \termino{GeoPackage} en \texttt{data/processed/}, 
todas proyectadas en \textsc{epsg}:32614 (UTM zona~14N), referencia cartográfica adecuada 
para la extensión de la \zmvm{}. Las capas se organizan en tres categorías temáticas, ilustradas en la 
Figura~\ref{fig:tg-capas}: (i)~\textbf{contexto territorial}, con los límites político-administrativos de la \cdmx{}, del 
Estado de México, el polígono de estudio y las variables de medio físico natural; (ii)~\textbf{red e infraestructura}, 
con las geometrías de los sistemas de transporte, la infraestructura vial y el trazado del Periférico; y 
(iii)~\textbf{indicadores del modelo VFT}, con las capas derivadas del cálculo de cobertura espacial, 
fuerza capilar y factor de desviación descritos en el Capítulo~\ref{ch:resultados-vft}.

\begin{figure}[htbp]
\centering
\begin{tikzpicture}[
  hdr1/.style = {rectangle, rounded corners=3pt,
                 draw=gray!60, fill=gray!55,
                 minimum width=4.2cm, minimum height=0.75cm,
                 align=center, font=\small\bfseries\color{white}},
  lyr1/.style = {rectangle, rounded corners=2pt,
                 draw=gray!45, fill=gray!8,
                 minimum width=4.2cm, minimum height=0.65cm,
                 align=center, font=\footnotesize},
  hdr2/.style = {rectangle, rounded corners=3pt,
                 draw=unamAzul, fill=unamAzul!70,
                 minimum width=4.2cm, minimum height=0.75cm,
                 align=center, font=\small\bfseries\color{white}},
  lyr2/.style = {rectangle, rounded corners=2pt,
                 draw=unamAzul!50, fill=unamAzul!8,
                 minimum width=4.2cm, minimum height=0.65cm,
                 align=center, font=\footnotesize},
  hdr3/.style = {rectangle, rounded corners=3pt,
                 draw=unamOro!80, fill=unamOro!75,
                 minimum width=4.2cm, minimum height=0.75cm,
                 align=center, font=\small\bfseries\color{white}},
  lyr3/.style = {rectangle, rounded corners=2pt,
                 draw=unamOro!60, fill=unamOro!10,
                 minimum width=4.2cm, minimum height=0.65cm,
                 align=center, font=\footnotesize},
  node distance = 0.12cm
]
%% — Columna 1: Contexto territorial —
\node[hdr1] (h1) at (0,0) {Contexto territorial};
\node[lyr1, below=of h1] (c1a) {\texttt{LimitesPoliticos.gpkg}};
\node[lyr1, below=of c1a] (c1b) {\texttt{PoligonosEdoMex.gpkg}};
\node[lyr1, below=of c1b] (c1c) {\texttt{PoligonoEstudio1.gpkg}};
\node[lyr1, below=of c1c] (c1d) {\texttt{MedioFisicoNatural.gpkg}};

%% — Columna 2: Red e infraestructura —
\node[hdr2] (h2) at (4.6,0) {Red e infraestructura};
\node[lyr2, below=of h2] (c2a) {\texttt{Transporte.gpkg}};
\node[lyr2, below=of c2a] (c2b) {\texttt{InfraEstructura.gpkg}};
\node[lyr2, below=of c2b] (c2c) {\texttt{Vialidades.gpkg}};

%% — Columna 3: Indicadores del modelo —
\node[hdr3] (h3) at (9.2,0) {Indicadores del modelo};
\node[lyr3, below=of h3] (c3a) {\texttt{Cobertura800m.gpkg}};
\node[lyr3, below=of c3a] (c3b) {\texttt{FuerzaCapilar.gpkg}};
\node[lyr3, below=of c3b] (c3c) {\texttt{FactorDesviacion.gpkg}};
\end{tikzpicture}
\caption[Organización de las capas GeoPackage de Transport-gis]%
  {Las diez capas \termino{GeoPackage} del directorio \texttt{data/processed/}, organizadas por 
  categoría temática. Todas están proyectadas en \textsc{epsg}:32614 (UTM zona~14N) y versionadas con Git~\textsc{lfs}.}
\label{fig:tg-capas}
\end{figure}

% =============================================================================
\subsubsection{Proyectos cartográficos y plantillas institucionales}
\label{subsubsec:tg-proyectos}
% =============================================================================

La Figura~\ref{fig:tg-arbol} muestra la estructura de directorios completa de \textbf{Transport-gis}. El repositorio se articula
en cinco carpetas de primer nivel con responsabilidades bien delimitadas: \texttt{analysis/} centraliza los \textit{scripts} de
procesamiento y los cuadernos de exploración; \texttt{data/processed/} almacena las diez capas \textit{GeoPackage} versionadas en
Git~\textsc{lfs}; \texttt{maps/} agrupa los proyectos QGIS~(\texttt{.qgz}), las plantillas visuales~(\texttt{.qpt}) con las paletas
institucionales de la \textsc{unam} y el directorio de exportaciones; \texttt{pipeline/} contiene el orquestador del proceso; y
\texttt{setup/} distribuye los activos tipográficos para garantizar la reproducibilidad visual en distintos sistemas operativos.
Los mapas terminados se exportan en \texttt{.pdf} (alta resolución para la tesis) y \texttt{.png} (presentaciones del coloquio),
cerrando el ciclo descrito en la Figura~\ref{fig:tg-pipeline}.

\begin{figure}[htbp]
\centering
\iffalse
\begin{forest}
  for tree = {
    font   = \footnotesize\ttfamily,
    grow'  = 0,
    child anchor  = west,
    parent anchor = east,
    anchor = west,
    draw   = gray!45,
    fill   = gray!7,
    rounded corners = 2pt,
    inner xsep = 6pt,
    inner ysep = 3pt,
    l sep  = 20pt,
    s sep  = 5pt,
    edge path = {
      \noexpand\path [draw=gray!50, thin, rounded corners=2pt]
        (!u.parent anchor) -- +(8pt,0) |- (.child anchor)
        \forestoption{edge label};
    },
  }
  [Transport-gis-zmvm-mjg/, fill=unamAzul!12, draw=unamAzul!70,
                             font=\small\bfseries\ttfamily
    [analysis/
      [notebooks/]
      [scripts/
        [apimetro\_client.py,
          draw=none, fill=none, font=\footnotesize\ttfamily\color{unamAzul}]
        [vft\_runner.py,
          draw=none, fill=none, font=\footnotesize\ttfamily\color{unamAzul}]
      ]
    ]
    [data/processed/
      [{\textit{10 capas .gpkg} \textnormal{(Git~LFS)}},
        draw=none, fill=none, font=\footnotesize\itshape\ttfamily]
    ]
    [maps/, fill=unamOro!8, draw=unamOro!60
      [exports/
        [{\textit{*.pdf · *.png}},
          draw=none, fill=none, font=\footnotesize\itshape\ttfamily]
      ]
      [fonts/]
      [projects/
        [\textit{*.qgz},
          draw=none, fill=none, font=\footnotesize\itshape\ttfamily]
      ]
      [templates/
        [{\textit{*.qpt} \textnormal{(paletas UNAM)}},
          draw=none, fill=none, font=\footnotesize\itshape\ttfamily]
      ]
    ]
    [pipeline/
      [run\_pipeline.py,
        draw=none, fill=none, font=\footnotesize\ttfamily\color{unamAzul}]
    ]
    [setup/
      [install\_fonts.sh,
        draw=none, fill=none, font=\footnotesize\ttfamily\color{unamAzul}]
      [install\_fonts.bat,
        draw=none, fill=none, font=\footnotesize\ttfamily\color{unamAzul}]
    ]
  ]
\end{forest}
\fi
\caption[Estructura de directorios de Transport-gis-zmvm-mjg]%
  {Árbol de directorios del repositorio \textbf{Transport-gis-zmvm-mjg}. Las carpetas \texttt{maps/} (oro) y la raíz (azul UNAM)
  son los nodos centrales de la producción cartográfica. Los archivos en azul son \textit{scripts} ejecutables;
  los de fondo plano son activos generados o de datos.}
\label{fig:tg-arbol}
\end{figure}

La Figura~\ref{fig:mapa-red-esquematica} presenta el resultado de este proceso aplicado a la red de transporte completa de la \zmvm{}:
un mapa de síntesis que integra los doce sistemas reconocidos por el modelo, proyectados en el espacio metropolitano real. Este mapa
constituye la representación cartográfica de referencia del estudio y es el punto de partida para la interpretación de los indicadores
presentados en el Capítulo~\ref{ch:resultados-vft}. Las versiones de alta resolución a página completa de todos los mapas producidos
por \textbf{Transport-gis} se encuentran en el Anexo~\ref{anx:cartografia-presentacion}.

\begin{figure}[htbp]
  \centering
  \figinclude[width=\textwidth]{Figures/Mapas/Mapa_Red_Esquematica}
  \caption[Red multimodal de transporte de la ZMVM]{Red esquemática multimodal de la \zmvm{} generada por \textbf{Transport-gis-zmvm-mjg}.
    Incluye los doce sistemas de transporte reconocidos por el modelo, proyectados en EPSG:32614 (UTM zona~14N) sobre los límites político-administrativos
    de la \cdmx{} y el Estado de México. \fuentefigura{Elaboración propia con QGIS~3.x, plantillas institucionales \textsc{unam};
    datos: \textsc{gtfs} \textsc{semovi} 2024 e \textsc{inegi} 2023.}}
  \label{fig:mapa-red-esquematica}
\end{figure}


%% ── 5-Resultados/5-2-1-Fase1/5-2-1-2-Cobertura.tex
\subsection*{5 2 1 2 Cobertura}
% !TeX root = ../../tesis.tex
% =============================================================================
%  5.2.1.2 NIVEL DE COBERTURA ESPACIAL
% =============================================================================
\subsection{\bajoRevision{Nivel de Cobertura Espacial (\texorpdfstring{$C$}{C})}}
\label{subsec:cobertura-espacial}
\notaRevision{Los resultados de esta subsección están sujetos a revisión con el
tutor. Los valores presentados corresponden a la corrida del modelo del
18 de abril de 2026 con un isócrono de 800~m y pueden ajustarse tras
validación metodológica.}

El \textbf{Nivel de Cobertura Espacial} ($C$) evalúa la equidad territorial de
la red de transporte masivo al cuantificar qué proporción del área de la
demarcación se encuentra dentro de un radio de caminata aceptable respecto
a alguna estación. Se define mediante la siguiente expresión:

\begin{equation}
    C = \frac{A_{\text{cubierta}}}{A_{\text{total}}} \times 100
    \label{eq:cobertura-res}
\end{equation}

donde $A_{\text{cubierta}}$ es el área resultante de la unión de los
isócronos peatonales (\termino{buffers} de 800~m) de todas las estaciones del sistema,
y $A_{\text{total}}$ es la superficie geodésica total de la demarcación
político-administrativa analizada.

El Cuadro~\ref{tab:peores-cdmx} identifica las cinco alcaldías de la
\textsc{cdmx} con menor cobertura del sistema de transporte masivo,
evidenciando las zonas de mayor exclusión territorial.

\begin{table}[H]
    \centering
    \caption[Alcaldías con menor cobertura de transporte masivo (CDMX)]{
    Cinco alcaldías de la Ciudad de México con menor porcentaje de cobertura
    territorial del sistema de transporte masivo. Isócrono: 800~m.
    Corrida: 2026-04-18.}
    \label{tab:peores-cdmx}
    \begin{tabular}{|l|r|r|r|}
        \hline
        \textbf{Demarcación} & \textbf{Área total (km²)} &
        \textbf{Área cubierta (km²)} & \textbf{Cobertura (\%)} \\
        \hline
        Milpa Alta              & 298.01 &  35.0    &  11.75 \\
        La Magdalena Contreras  &  63.37 &  16.33   &  25.76 \\
        Tlalpan                 & 314.25 &  94.31   &  30.01 \\
        Tláhuac                 &  85.78 &  43.85   &  51.12 \\
        Cuajimalpa de Morelos   &  71.10 &  37.12   &  52.20 \\
        \hline
    \end{tabular}
    \fuentefigura{Elaboración propia con \texttt{VFTModel} y datos de
    \texttt{Apimetro}.}
\end{table}

\begin{table}[H]
    \centering
    \caption[Municipios del EDOMEX con menor cobertura de transporte masivo]{
    Cinco municipios del Estado de México con menor porcentaje de cobertura
    territorial del sistema de transporte masivo. Isócrono: 800~m.
    Corrida: 2026-04-18.}
    \label{tab:peores-edomex}
    \begin{tabular}{|l|r|r|r|}
        \hline
        \textbf{Demarcación} & \textbf{Área total (km²)} &
        \textbf{Área cubierta (km²)} & \textbf{Cobertura (\%)} \\
        \hline
        Tonanitla       &   8.52 & 0.00 & 0.00 \\
        Capulhuac       &  21.50 & 0.00 & 0.00 \\
        Calimaya        & 104.26 & 0.00 & 0.00 \\
        Donato Guerra   & 181.36 & 0.00 & 0.00 \\
        Ayapango        &  50.83 & 0.00 & 0.00 \\
        \hline
    \end{tabular}
    \fuentefigura{Elaboración propia con \texttt{VFTModel} y datos de
    \texttt{Apimetro}.}
\end{table}

% =============================================================================
%  Desglose de Cobertura por Sistema de Transporte
% =============================================================================

\subsubsection{Sistema de Transporte Colectivo Metro (METRO)}
\label{subsubsec:cobertura-metro}

El Sistema de Transporte Colectivo Metro constituye la columna vertebral de la
movilidad masiva en la Ciudad de México. Debido a su naturaleza de red ferroviaria
subterránea y elevada con derecho de vía exclusivo, sus estaciones actúan como
nodos estructurales de alta capacidad que articulan los flujos intermodales de la
metrópoli. El análisis de su cobertura espacial mediante isócronos peatonales de
800~m permite identificar las demarcaciones que se benefician directamente de este
sistema, así como aquellas periferias y zonas centrales que presentan una
desconexión crítica respecto a la red de transporte más masiva de la entidad.

El Cuadro~\ref{tab:cobertura-metro-cdmx} detalla las cinco alcaldías de la Ciudad de
México con los niveles más altos de cobertura territorial respecto a la red del
Metro, evidenciando la concentración de la infraestructura ferroviaria en el núcleo
urbano denso de la cuenca.

\begin{table}[H]
    \centering
    \caption[Alcaldías con mayor cobertura del STC Metro (CDMX)]{
    Cinco alcaldías de la Ciudad de México con mayor porcentaje de cobertura
    territorial para el Sistema de Transporte Colectivo Metro. Isócrono: 800~m.
    Corrida: 2026-04-18.}
    \label{tab:cobertura-metro-cdmx}
    \begin{tabular}{|l|r|r|r|}
        \hline
        \textbf{Demarcación} & \textbf{Área total (km²)} & \textbf{Área cubierta (km²)} & \textbf{Cobertura (\%)} \\
        \hline
        Cuauhtémoc          & 32.50 &  28.35 &  87.23 \\
        Benito Juárez       & 26.68 &  21.34 &  79.99 \\
        Venustiano Carranza & 33.84 &  24.39 &  72.77 \\
        Iztacalco           & 23.08 &  11.66 &  50.52 \\
        Azcapotzalco        & 33.50 &  13.80 &  41.18 \\
        \hline
    \end{tabular}
    \fuentefigura{Elaboración propia con \texttt{VFTModel} y datos de \texttt{Apimetro}.}
\end{table}

Por su parte, la distribución geográfica de esta cobertura se visualiza en la
Figura~\ref{fig:mapa-cobertura-metro}, donde se aprecia la concentración de la
infraestructura en la zona central de la cuenca y el drástico desvanecimiento del
servicio hacia las demarcaciones del sur y del poniente.

\begin{figure}[H]
    \centering
    \figinclude[width=0.85\textwidth]{Figures/Cap5/Cobertura/Mapas/NB02_fig02_METRO_cobertura.png}
    \caption[Mapa de cobertura territorial del STC Metro en la CDMX]{Mapa del
    nivel de cobertura territorial del Sistema de Transporte Colectivo Metro
    por demarcación en la Ciudad de México, calculado a partir de \termino{buffers}
    de caminata peatonal de 800~m.}
    \label{fig:mapa-cobertura-metro}
    \fuentefigura{Elaboración propia con \texttt{VFTModel} utilizando visualización geoespacial.}
\end{figure}

\subsubsection{Metrobús (MB)}
\label{subsubsec:cobertura-mb}

El sistema Metrobús, que opera bajo el modelo de Autobuses de Tránsito Rápido
(\termino{Bus Rapid Transit}, BRT), se ha consolidado como la segunda red de
transporte masivo más importante de la Ciudad de México. Dada su mayor flexibilidad
constructiva frente a los sistemas ferroviarios pesados, su trazado ha logrado
expandirse por corredores viales estratégicos, complementando y alimentando a la
red del Metro. No obstante, al evaluar su cobertura mediante isócronos peatonales
de 800~m, persisten notables brechas territoriales, dejando a varias demarcaciones
periféricas con acceso limitado o nulo a esta red troncal.

El Cuadro~\ref{tab:cobertura-mb-cdmx} presenta las cinco alcaldías con mayor
proporción de su territorio cubierto por el Metrobús, destacando la concentración
del servicio en las zonas centro, norte y oriente de la ciudad, siguiendo los ejes
viales primarios.

\begin{table}[H]
    \centering
    \caption[Alcaldías con mayor cobertura de Metrobús (CDMX)]{
    Cinco alcaldías de la Ciudad de México con mayor porcentaje de cobertura
    territorial para el sistema Metrobús (MB). Isócrono: 800~m.
    Corrida: 2026-04-18.}
    \label{tab:cobertura-mb-cdmx}
    \begin{tabular}{|l|r|r|r|}
        \hline
        \textbf{Demarcación} & \textbf{Área total (km²)} & \textbf{Área cubierta (km²)} & \textbf{Cobertura (\%)} \\
        \hline
        Cuauhtémoc          & 32.50 &  26.74 &  82.27 \\
        Benito Juárez       & 26.68 &  16.21 &  60.74 \\
        Iztacalco           & 23.08 &  13.95 &  60.44 \\
        Venustiano Carranza & 33.84 &  16.10 &  47.58 \\
        Gustavo A. Madero   & 87.84 &  38.75 &  44.12 \\
        \hline
    \end{tabular}
    \fuentefigura{Elaboración propia con \texttt{VFTModel} y datos de \texttt{Apimetro}.}
\end{table}

Visualmente, el alcance espacial del sistema se ilustra en la
Figura~\ref{fig:mapa-cobertura-mb}. El mapa confirma cómo las líneas de BRT se
concentran fuertemente en las zonas centro, norte y oriente —siguiendo los ejes
viales primarios—, mientras que las barreras topográficas y de traza urbana
mantienen a las demarcaciones del surponiente con una cobertura residual o
inexistente.

\begin{figure}[H]
    \centering
    \figinclude[width=0.85\textwidth]{Figures/Cap5/Cobertura/Mapas/NB02_fig02_MB_cobertura.png}
    \caption[Mapa de cobertura territorial del Metrobús en la CDMX]{Mapa del
    nivel de cobertura territorial del sistema Metrobús por demarcación en la
    Ciudad de México, calculado a partir de \termino{buffers} de caminata
    peatonal de 800~m.}
    \label{fig:mapa-cobertura-mb}
    \fuentefigura{Elaboración propia con \texttt{VFTModel} utilizando visualización geoespacial.}
\end{figure}

\subsubsection{Cablebús (CBB)}
\label{subsubsec:cobertura-cbb}

El Cablebús \DIFdelbegin \DIFdel{representa un paradigma distinto dentro de la red de transporte masivo.
Al ser un }\DIFdelend \DIFaddbegin \DIFadd{es un }\DIFaddend sistema de teleférico urbano \DIFdelbegin \DIFdel{, su diseño no busca abarcar grandes
extensiones del territorio mediante una red interconectada, sino }\DIFdelend \DIFaddbegin \DIFadd{diseñado para }\DIFaddend resolver problemas de
accesibilidad \DIFdelbegin \DIFdel{muy puntuales }\DIFdelend en zonas de alta marginación y topografía compleja\DIFdelbegin \DIFdel{(principalmente en la periferia). Por lo tanto, su nivel de cobertura espacial a
escala metropolitana es inherentemente focalizado, }\DIFdelend \DIFaddbegin \DIFadd{, }\DIFaddend actuando como
alimentador de alta capacidad hacia \DIFdelbegin \DIFdel{las líneas del Metro. }%DIFDELCMD < 

%DIFDELCMD < %%%
\DIFdel{El Cuadro~\ref{tab:cobertura-cbb-cdmx} presenta las cinco alcaldías con mayor
cobertura del Cablebús. Como es de esperarse por la naturaleza del sistema, la gran
mayoría de las demarcaciones de la Ciudad de México carecen por completo de esta
infraestructura, por lo que incluso los valores más altos son reducidos en términos
porcentuales absolutos}\DIFdelend \DIFaddbegin \DIFadd{el Metro. Su cobertura se concentra en
Gustavo A. Madero (12.88\,\%), Miguel Hidalgo (11.85\,\%) e Iztapalapa
(11.27\,\%), con presencia marginal en el resto de la ciudad
(Cuadro~\ref{tab:cobertura-cbb-cdmx})}\DIFaddend .

\begin{table}[H]
    \centering
    \caption[Alcaldías con mayor cobertura de Cablebús (CDMX)]{
    Cinco alcaldías de la Ciudad de México con mayor porcentaje de cobertura
    territorial para el sistema Cablebús (CBB). Isócrono: 800~m.
    Corrida: 2026-04-18.}
    \label{tab:cobertura-cbb-cdmx}
    \begin{tabular}{|l|r|r|r|}
        \hline
        \textbf{Demarcación} & \textbf{Área total (km²)} & \textbf{Área cubierta (km²)} & \textbf{Cobertura (\%)} \\
        \hline
        Gustavo A. Madero   &  87.84 & 11.31 &  12.88 \\
        Miguel Hidalgo      &  46.39 &  5.50 &  11.85 \\
        Iztapalapa          & 113.07 & 12.74 &  11.27 \\
        Álvaro Obregón      &  95.82 &  4.38 &   4.57 \\
        Venustiano Carranza &  33.84 &  0.00 &   0.00 \\
        \hline
    \end{tabular}
    \fuentefigura{Elaboración propia con \texttt{VFTModel} y datos de \texttt{Apimetro}.}
\end{table}

La Figura~\ref{fig:mapa-cobertura-cbb} ilustra esta focalización espacial. El mapa
revela cómo las áreas de influencia del Cablebús se restringen a corredores muy
específicos (como en Gustavo A. Madero, Iztapalapa y Álvaro Obregón), contrastando
drásticamente con el vacío de infraestructura aérea en el resto de la cuenca.

\begin{figure}[H]
    \centering
    \figinclude[width=0.85\textwidth]{Figures/Cap5/Cobertura/Mapas/NB02_fig02_CBB_cobertura.png}
    \caption[Mapa de cobertura territorial del Cablebús en la CDMX]{Mapa del
    nivel de cobertura territorial del sistema Cablebús por demarcación en la
    Ciudad de México, calculado a partir de \termino{buffers} de caminata
    peatonal de 800~m.}
    \label{fig:mapa-cobertura-cbb}
    \fuentefigura{Elaboración propia con \texttt{VFTModel} utilizando visualización geoespacial.}
\end{figure}

\subsubsection{Red de Transporte de Pasajeros (RTP)}
\label{subsubsec:cobertura-rtp}

La Red de Transporte de Pasajeros (\textsc{rtp}) desempeña un rol fundamental en
la equidad espacial del sistema de movilidad de la Ciudad de México. A diferencia
de los sistemas confinados, la \textsc{rtp} opera como un servicio de superficie
con vocación social y de alimentación, diseñado específicamente para penetrar en
zonas periféricas, de topografía compleja y de menor densidad, donde la
implementación de infraestructura pesada es inviable. Su extensión territorial es
notablemente mayor que la de cualquier otro sistema de la metrópoli.

El Cuadro~\ref{tab:cobertura-rtp-cdmx} presenta las cinco alcaldías con mayor
proporción de cobertura por parte de la \textsc{rtp}, confirmando la concentración
del servicio en el centro y norte de la ciudad, donde la densidad de rutas y paradas
alcanza sus niveles más altos.

\begin{table}[H]
    \centering
    \caption[Alcaldías con mayor cobertura de RTP (CDMX)]{
    Cinco alcaldías de la Ciudad de México con mayor porcentaje de cobertura
    territorial para la Red de Transporte de Pasajeros (RTP). Isócrono: 800~m.
    Corrida: 2026-04-18.}
    \label{tab:cobertura-rtp-cdmx}
    \begin{tabular}{|l|r|r|r|}
        \hline
        \textbf{Demarcación} & \textbf{Área total (km²)} & \textbf{Área cubierta (km²)} & \textbf{Cobertura (\%)} \\
        \hline
        Miguel Hidalgo      &  46.39 &  41.90 &  90.32 \\
        Cuauhtémoc          &  32.50 &  29.32 &  90.22 \\
        Azcapotzalco        &  33.50 &  30.08 &  89.81 \\
        Gustavo A. Madero   &  87.84 &  75.86 &  86.37 \\
        Iztapalapa          & 113.07 &  93.66 &  82.83 \\
        \hline
    \end{tabular}
    \fuentefigura{Elaboración propia con \texttt{VFTModel} y datos de \texttt{Apimetro}.}
\end{table}

La Figura~\ref{fig:mapa-cobertura-rtp} muestra la dispersión territorial de este
sistema. A diferencia del Metro o del Metrobús, el mapa de la \textsc{rtp} revela
una capilaridad mucho mayor hacia el sur y poniente de la ciudad, confirmando su
rol estratégico como conector de las zonas más alejadas de la cuenca con la red
troncal.

\begin{figure}[H]
    \centering
    \figinclude[width=0.85\textwidth]{Figures/Cap5/Cobertura/Mapas/NB02_fig02_RTP_cobertura.png}
    \caption[Mapa de cobertura territorial de la RTP en la CDMX]{Mapa del
    nivel de cobertura territorial de la Red de Transporte de Pasajeros
    (\textsc{rtp}) por demarcación en la Ciudad de México, calculado a partir
    de \termino{buffers} de caminata peatonal de 800~m.}
    \label{fig:mapa-cobertura-rtp}
    \fuentefigura{Elaboración propia con \texttt{VFTModel} utilizando visualización geoespacial.}
\end{figure}

\subsubsection{Trolebús (TROLE)}
\label{subsubsec:cobertura-trole}

El \DIFaddbegin \DIFadd{Trolebús —operado por el }\DIFaddend Servicio de Transportes Eléctricos (\textsc{ste})\DIFdelbegin \DIFdel{, a través de su red de
Trolebús, constituye }\DIFdelend \DIFaddbegin \DIFadd{—
es }\DIFaddend uno de los ejes históricos \DIFdelbegin \DIFdel{y fundamentales de la
}\DIFdelend \DIFaddbegin \DIFadd{de }\DIFaddend electromovilidad de cero emisiones en la capital\DIFdelbegin \DIFdel{. Su red se distribuye
estratégicamente sobre los principales ejes viales,
concentrando su mayor área
de influencia en las zonas centro , norte }\DIFdelend \DIFaddbegin \DIFadd{,
con presencia relevante en el centro }\DIFaddend y oriente de la ciudad. \DIFaddbegin \DIFadd{El
Cuadro~\ref{tab:cobertura-trole-cdmx} registra coberturas superiores al 50\,\%
en cuatro alcaldías centrales, encabezadas por Cuauhtémoc (69.90\,\%).
}\DIFaddend 

\textbf{Nota metodológica:} Es importante precisar que, al momento de la
ejecución de este análisis espacial (18 de abril de 2026), el Portal de Datos
Abiertos de la Ciudad de México aún no había liberado los trazados geoespaciales
ni el padrón de paradas oficiales correspondientes a la \textbf{Línea 13} del
Trolebús. Por consiguiente, dicha línea no se encuentra contabilizada en los
isócronos de cobertura presentados en esta sección.

\DIFdelbegin \DIFdel{El Cuadro~\ref{tab:cobertura-trole-cdmx} enumera las cinco alcaldías con mayor
cobertura territorial de este sistema, todas ellas ubicadas en la zona central y
oriente de la ciudad, donde se concentra la red eléctrica de superficie.
}%DIFDELCMD < 

%DIFDELCMD < %%%
\DIFdelend \begin{table}[H]
    \centering
    \caption[Alcaldías con mayor cobertura de Trolebús (CDMX)]{
    Cinco alcaldías de la Ciudad de México con mayor porcentaje de cobertura
    territorial para el sistema Trolebús (TROLE). Isócrono: 800~m.
    Corrida: 2026-04-18.}
    \label{tab:cobertura-trole-cdmx}
    \begin{tabular}{|l|r|r|r|}
        \hline
        \textbf{Demarcación} & \textbf{Área total (km²)} & \textbf{Área cubierta (km²)} & \textbf{Cobertura (\%)} \\
        \hline
        Cuauhtémoc    &  32.50 & 22.72 &  69.90 \\
        Benito Juárez &  26.68 & 16.50 &  61.85 \\
        Azcapotzalco  &  33.50 & 17.48 &  52.18 \\
        Iztacalco     &  23.08 & 11.93 &  51.69 \\
        Coyoacán      & 125.17 & 25.17 &  46.71 \\
        \hline
    \end{tabular}
    \fuentefigura{Elaboración propia con \texttt{VFTModel} y datos de \texttt{Apimetro}.}
\end{table}

Por su parte, la Figura~\ref{fig:mapa-cobertura-trole} ilustra la huella espacial
del Trolebús. El mapa visibiliza la fuerte concentración del servicio en alcaldías
como Benito Juárez, Cuauhtémoc, Iztacalco e Iztapalapa (impulsada por el Trolebús
Elevado), contrastando drásticamente con la carencia de infraestructura eléctrica
hacia las periferias montañosas.

\begin{figure}[H]
    \centering
    \figinclude[width=0.85\textwidth]{Figures/Cap5/Cobertura/Mapas/NB02_fig02_TROLE_cobertura.png}
    \caption[Mapa de cobertura territorial del Trolebús en la CDMX]{Mapa del
    nivel de cobertura territorial del sistema Trolebús por demarcación en la
    Ciudad de México, calculado a partir de \termino{buffers} de caminata
    peatonal de 800~m.}
    \label{fig:mapa-cobertura-trole}
    \fuentefigura{Elaboración propia con \texttt{VFTModel} utilizando visualización geoespacial.}
\end{figure}

\subsubsection{Corredores Concesionados (CC)}
\label{subsubsec:cobertura-cc}

Los Corredores Concesionados constituyen la red de rutas de transporte de
superficie operadas por empresas privadas con concesión del Gobierno de la Ciudad
de México. Con base en su lógica de trazado sobre los ejes viales primarios y
secundarios de la ciudad consolidada, este sistema presenta uno de los niveles de
cobertura territorial más altos de la metrópoli, especialmente en las alcaldías de
mayor densidad urbana del centro y el oriente. Más que un sistema residual, los
Corredores Concesionados actúan como la red capilar que penetra en los intersticios
que los sistemas troncales no pueden servir directamente.

El Cuadro~\ref{tab:cobertura-cc-cdmx} presenta las cinco alcaldías con mayor
cobertura territorial de este sistema, donde se aprecia la saturación práctica del
servicio en el núcleo urbano denso de la capital.

\begin{table}[H]
    \centering
    \caption[Alcaldías con mayor cobertura de Corredores Concesionados (CDMX)]{
    Cinco alcaldías de la Ciudad de México con mayor porcentaje de cobertura
    territorial para el sistema de Corredores Concesionados (CC). Isócrono: 800~m.
    Corrida: 2026-04-18.}
    \label{tab:cobertura-cc-cdmx}
    \begin{tabular}{|l|r|r|r|}
        \hline
        \textbf{Demarcación} & \textbf{Área total (km²)} & \textbf{Área cubierta (km²)} & \textbf{Cobertura (\%)} \\
        \hline
        Cuauhtémoc          & 32.50 & 32.47 &  99.90 \\
        Benito Juárez       & 26.68 & 26.32 &  98.65 \\
        Iztacalco           & 23.08 & 22.25 &  96.42 \\
        Miguel Hidalgo      & 46.39 & 44.33 &  95.57 \\
        Venustiano Carranza & 33.84 & 27.77 &  82.06 \\
        \hline
    \end{tabular}
    \fuentefigura{Elaboración propia con \texttt{VFTModel} y datos de \texttt{Apimetro}.}
\end{table}

La Figura~\ref{fig:mapa-cobertura-cc} confirma la distribución omnipresente de
este sistema en la ciudad consolidada. El mapa ilustra cómo los Corredores
Concesionados recubren de manera prácticamente total las alcaldías centrales,
mientras que el servicio se va diluyendo hacia las periferias con suelo de
conservación o menor densidad.

\begin{figure}[H]
    \centering
    \figinclude[width=0.85\textwidth]{Figures/Cap5/Cobertura/Mapas/NB02_fig02_CC_cobertura.png}
    \caption[Mapa de cobertura territorial de Corredores Concesionados en la CDMX]{
    Mapa del nivel de cobertura territorial del sistema de Corredores
    Concesionados por demarcación en la Ciudad de México, calculado a partir de
    \termino{buffers} de caminata peatonal de 800~m.}
    \label{fig:mapa-cobertura-cc}
    \fuentefigura{Elaboración propia con \texttt{VFTModel} utilizando visualización geoespacial.}
\end{figure}

\subsubsection{Tren Ligero (TL)}
\label{subsubsec:cobertura-tl}

El Tren Ligero es un sistema de transporte férreo de capacidad media operado por
el Servicio de Transportes Eléctricos (\textsc{ste}). Su trazado consta de una
única línea que conecta la estación terminal Tasqueña de la Línea 2 del Metro con
el centro de la alcaldía Xochimilco. Debido a esta configuración lineal y su
confinamiento exclusivo en el sur de la cuenca, su área de influencia geográfica
es altamente focalizada.

El Cuadro~\ref{tab:cobertura-tl-cdmx} presenta las alcaldías con mayor nivel de
cobertura por parte del Tren Ligero. Únicamente Coyoacán, Xochimilco y Tlalpan
registran acceso al sistema dentro del isócrono peatonal de 800~m, mientras que
el resto de las demarcaciones carece por completo de este servicio.

\begin{table}[H]
    \centering
    \caption[Alcaldías con mayor cobertura de Tren Ligero (CDMX)]{
    Cinco alcaldías de la Ciudad de México con mayor porcentaje de cobertura
    territorial para el sistema Tren Ligero (TL). Isócrono: 800~m.
    Corrida: 2026-04-18.}
    \label{tab:cobertura-tl-cdmx}
    \begin{tabular}{|l|r|r|r|}
        \hline
        \textbf{Demarcación} & \textbf{Área total (km²)} & \textbf{Área cubierta (km²)} & \textbf{Cobertura (\%)} \\
        \hline
        Coyoacán            &  53.88 & 9.31 &  17.28 \\
        Xochimilco          & 114.03 & 6.82 &   5.98 \\
        Tlalpan             & 314.25 & 4.07 &   1.30 \\
        Venustiano Carranza &  33.84 & 0.00 &   0.00 \\
        Tláhuac             &  85.78 & 0.00 &   0.00 \\
        \hline
    \end{tabular}
    \fuentefigura{Elaboración propia con \texttt{VFTModel} y datos de \texttt{Apimetro}.}
\end{table}

La Figura~\ref{fig:mapa-cobertura-tl} ilustra espacialmente esta restricción
territorial. El mapa confirma cómo el servicio es exclusivo del corredor sur,
actuando principalmente como alimentador de la red masiva central para los
habitantes de las alcaldías sureñas.

\begin{figure}[H]
    \centering
    \figinclude[width=0.85\textwidth]{Figures/Cap5/Cobertura/Mapas/NB02_fig02_TL_cobertura.png}
    \caption[Mapa de cobertura territorial del Tren Ligero en la CDMX]{Mapa del
    nivel de cobertura territorial del sistema Tren Ligero por demarcación en
    la Ciudad de México, calculado a partir de \termino{buffers} de caminata
    peatonal de 800~m.}
    \label{fig:mapa-cobertura-tl}
    \fuentefigura{Elaboración propia con \texttt{VFTModel} utilizando visualización geoespacial.}
\end{figure}

\subsubsection{Mexibús (MXB)}
\label{subsubsec:cobertura-mxb}

El sistema Mexibús \DIFdelbegin \DIFdel{es el equivalente al Metrobús en el Estado de México: una red
de autobuses de tránsito rápido (}%DIFDELCMD < \termino{Bus Rapid Transit}%%%
\DIFdel{, BRT) que }\DIFdelend opera sobre corredores viales del \DIFdelbegin \DIFdel{ámbito metropolitano conurbado. Su presencia }\DIFdelend \DIFaddbegin \DIFadd{Estado de México con una
presencia marginal en la Ciudad de México. Su cobertura máxima }\DIFaddend dentro de los
límites \DIFdelbegin \DIFdel{político-administrativos de la Ciudad de México es marginal, circunscrita
únicamente a los puntos de cruce o de conexión interestatal, por lo que su
cobertura territorial en las
alcaldías de la capital es forzosamente residual.
}%DIFDELCMD < 

%DIFDELCMD < %%%
\DIFdel{El }\DIFdelend \DIFaddbegin \DIFadd{capitalinos no supera el 6.58\,\% (Iztacalco), siendo nula en las
demarcaciones del sur y oriente (}\DIFaddend Cuadro~\ref{tab:cobertura-mxb-cdmx}\DIFdelbegin \DIFdel{identifica las alcaldías de la Ciudad de México con mayor cobertura territorial del sistema Mexibús. Los valores reducidos
—ninguno superior al 7\%— confirman que este sistema opera fundamentalmente fuera
de los límites de la entidad y que su incidencia en la cobertura intraurbana
capitalina es mínima}\DIFdelend \DIFaddbegin \DIFadd{). Para
la visualización cartográfica de este sistema en el contexto de la red
metropolitana véase la Figura~\ref{fig:anx-mapa-cobertura-total} del
Anexo~\ref{anx:cartografia-presentacion}}\DIFaddend .

\begin{table}[H]
    \centering
    \caption[Alcaldías con mayor cobertura de Mexibús (CDMX)]{
    Cinco alcaldías de la Ciudad de México con mayor porcentaje de cobertura
    territorial para el sistema Mexibús (MXB). Isócrono: 800~m.
    Corrida: 2026-04-18.}
    \label{tab:cobertura-mxb-cdmx}
    \begin{tabular}{|l|r|r|r|}
        \hline
        \textbf{Demarcación} & \textbf{Área total (km²)} & \textbf{Área cubierta (km²)} & \textbf{Cobertura (\%)} \\
        \hline
        Iztacalco           &  23.08 & 1.52 & 6.58 \\
        Venustiano Carranza &  33.84 & 1.90 & 5.62 \\
        Gustavo A. Madero   &  87.84 & 3.23 & 3.67 \\
        Xochimilco          & 114.03 & 0.00 & 0.00 \\
        Tláhuac             &  85.78 & 0.00 & 0.00 \\
        \hline
    \end{tabular}
    \fuentefigura{Elaboración propia con \texttt{VFTModel} y datos de \texttt{Apimetro}.}
\end{table}

\DIFdelbegin \DIFdel{La Figura~\ref{fig:mapa-cobertura-mxb} confirma esta presencia puntual del
sistema dentro de la capital. El mapa ilustra cómo la única área de influencia
del Mexibús en la Ciudad de México corresponde a los cruces fronterizos del
nororiente, sin penetrar en el tejido urbano interior.
}%DIFDELCMD < 

%DIFDELCMD < \begin{figure}[H]
%DIFDELCMD <     \centering
%DIFDELCMD <     \figinclude[width=0.85\textwidth]{Figures/Cap5/Cobertura/Mapas/NB02_fig02_MEXIBÚS_cobertura.png}
%DIFDELCMD <     %%%
%DIFDELCMD < \caption[Mapa de cobertura territorial de Mexibús en la CDMX]{%
{%DIFAUXCMD
\DIFdelFL{Mapa del
    nivel de cobertura territorial del sistema Mexibús por demarcación en la
    Ciudad de México, calculado a partir de }%DIFDELCMD < \termino{buffers} %%%
\DIFdelFL{de caminata
    peatonal de 800~m.}}
    %DIFAUXCMD
%DIFDELCMD < \label{fig:mapa-cobertura-mxb}
%DIFDELCMD <     \fuentefigura{Elaboración propia con \texttt{VFTModel} utilizando visualización geoespacial.}
%DIFDELCMD < \end{figure}
%DIFDELCMD < 

%DIFDELCMD < %%%
\DIFdelend \subsubsection{Mexicable (MEXICABLE)}
\label{subsubsec:cobertura-mexicable}

El Mexicable es un \DIFdelbegin \DIFdel{sistema de }\DIFdelend teleférico urbano \DIFdelbegin \DIFdel{que opera principalmente en el
}\DIFdelend \DIFaddbegin \DIFadd{del }\DIFaddend Estado de México \DIFdelbegin \DIFdel{, diseñado para sortear las barreras topográficas de las zonas
altas y densamente pobladas de la periferia metropolitana. Su }\DIFdelend \DIFaddbegin \DIFadd{cuya }\DIFaddend incursión en la
Ciudad de México es mínima\DIFdelbegin \DIFdel{y de carácter estrictamente interconectivo, fungiendo
como enlace alimentador hacia las redes troncales en la zona norte de la capital.
}%DIFDELCMD < 

%DIFDELCMD < %%%
\DIFdel{El Cuadro~\ref{tab:cobertura-mexicable-cdmx} presenta las alcaldías de la capital
con mayor nivel de cobertura para el Mexicable. Como es natural dada su jurisdicción
y trazado, }\DIFdelend \DIFaddbegin \DIFadd{: }\DIFaddend únicamente \DIFdelbegin \DIFdel{la alcaldía }\DIFdelend Gustavo A. Madero registra \DIFdelbegin \DIFdel{un área de influencia
efectiva dentro del isócrono peatonal de 800~m, mientras que el resto de las
demarcaciones presentan cobertura nula}\DIFdelend \DIFaddbegin \DIFadd{cobertura
efectiva (3.19\,\%), mientras el resto del territorio capitalino presenta valores
nulos (Cuadro~\ref{tab:cobertura-mexicable-cdmx}). Para la visualización
cartográfica de este sistema en el contexto de la red metropolitana véase la
Figura~\ref{fig:anx-mapa-cobertura-total} del
Anexo~\ref{anx:cartografia-presentacion}}\DIFaddend .

\begin{table}[H]
    \centering
    \caption[Alcaldías con mayor cobertura de Mexicable (CDMX)]{
    Cinco alcaldías de la Ciudad de México con mayor porcentaje de cobertura
    territorial para el sistema Mexicable. Isócrono: 800~m.
    Corrida: 2026-04-18.}
    \label{tab:cobertura-mexicable-cdmx}
    \begin{tabular}{|l|r|r|r|}
        \hline
        \textbf{Demarcación} & \textbf{Área total (km²)} & \textbf{Área cubierta (km²)} & \textbf{Cobertura (\%)} \\
        \hline
        Gustavo A. Madero   &  87.84 & 2.80 & 3.19 \\
        Venustiano Carranza &  33.84 & 0.00 & 0.00 \\
        Xochimilco          & 114.03 & 0.00 & 0.00 \\
        Tláhuac             &  85.78 & 0.00 & 0.00 \\
        Tlalpan             & 314.25 & 0.00 & 0.00 \\
        \hline
    \end{tabular}
    \fuentefigura{Elaboración propia con \texttt{VFTModel} y datos de \texttt{Apimetro}.}
\end{table}

\DIFdelbegin \DIFdel{La Figura~\ref{fig:mapa-cobertura-mexicable} visualiza la focalización extrema de
este servicio. El mapa ilustra claramente cómo la infraestructura aérea del
Mexicable se restringe a un pequeño polígono en el extremo norte de la ciudad,
reiterando su naturaleza como transporte de penetración desde la zona conurbada.
}%DIFDELCMD < 

%DIFDELCMD < \begin{figure}[H]
%DIFDELCMD <     \centering
%DIFDELCMD <     \figinclude[width=0.85\textwidth]{Figures/Cap5/Cobertura/Mapas/NB02_fig02_MEXICABLE_cobertura.png}
%DIFDELCMD <     %%%
%DIFDELCMD < \caption[Mapa de cobertura territorial del Mexicable en la CDMX]{%
{%DIFAUXCMD
\DIFdelFL{Mapa del
    nivel de cobertura territorial del sistema Mexicable por demarcación en la
    Ciudad de México, calculado a partir de }%DIFDELCMD < \termino{buffers} %%%
\DIFdelFL{de caminata
    peatonal de 800~m.}}
    %DIFAUXCMD
%DIFDELCMD < \label{fig:mapa-cobertura-mexicable}
%DIFDELCMD <     \fuentefigura{Elaboración propia con \texttt{VFTModel} utilizando visualización geoespacial.}
%DIFDELCMD < \end{figure}
%DIFDELCMD < 

%DIFDELCMD < %%%
\DIFdelend \subsubsection{Tren Suburbano (SUB)}
\label{subsubsec:cobertura-sub}

El Tren Suburbano \DIFdelbegin \DIFdel{de la }%DIFDELCMD < \zmvm{} %%%
\DIFdel{es un sistema ferroviario regional de alta
capacidad que }\DIFdelend conecta la terminal Buenavista \DIFdelbegin \DIFdel{, en la alcaldía Cuauhtémoc, con el
municipio de }\DIFdelend \DIFaddbegin \DIFadd{con }\DIFaddend Cuautitlán \DIFdelbegin \DIFdel{, en el }\DIFdelend \DIFaddbegin \DIFadd{(}\DIFaddend Estado de
México\DIFdelbegin \DIFdel{. Su trazado, mayoritariamente
confinado al corredor norte de la cuenca metropolitana, limita }\DIFdelend \DIFaddbegin \DIFadd{), limitando }\DIFaddend su área de influencia \DIFdelbegin \DIFdel{directa dentro de }\DIFdelend \DIFaddbegin \DIFadd{en }\DIFaddend la Ciudad de México \DIFdelbegin \DIFdel{a un corredor muy acotado.
Para
las alcaldías del sur y oriente, este sistema resulta prácticamente inaccesible
sin una transferencia modal significativa.
}%DIFDELCMD < 

%DIFDELCMD < %%%
\DIFdel{El }\DIFdelend \DIFaddbegin \DIFadd{al corredor norte.
Solo Cuauhtémoc (6.18\,\%) y Azcapotzalco (5.99\,\%) registran cobertura
efectiva, con el resto del territorio capitalino sin acceso directo al sistema
(}\DIFaddend Cuadro~\ref{tab:cobertura-sub-cdmx}\DIFdelbegin \DIFdel{detalla las cinco alcaldías de la Ciudad de
México con mayor nivel de cobertura por parte del Tren Suburbano, evidenciando la
concentración casi exclusiva del servicio en las demarcaciones del norte y centro
de la capital}\DIFdelend \DIFaddbegin \DIFadd{)}\DIFaddend .

\begin{table}[H]
    \centering
    \caption[Alcaldías con mayor cobertura de Tren Suburbano (CDMX)]{
    Cinco alcaldías de la Ciudad de México con mayor porcentaje de cobertura
    territorial para el sistema Tren Suburbano (SUB). Isócrono: 800~m.
    Corrida: 2026-04-18.}
    \label{tab:cobertura-sub-cdmx}
    \begin{tabular}{|l|r|r|r|}
        \hline
        \textbf{Demarcación} & \textbf{Área total (km²)} & \textbf{Área cubierta (km²)} & \textbf{Cobertura (\%)} \\
        \hline
        Cuauhtémoc          &  32.50 & 2.01 & 6.18 \\
        Azcapotzalco        &  33.50 & 2.01 & 5.99 \\
        Xochimilco          & 114.03 & 0.00 & 0.00 \\
        Venustiano Carranza &  33.84 & 0.00 & 0.00 \\
        Tláhuac             &  85.78 & 0.00 & 0.00 \\
        \hline
    \end{tabular}
    \fuentefigura{Elaboración propia con \texttt{VFTModel} y datos de \texttt{Apimetro}.}
\end{table}

La Figura~\ref{fig:mapa-cobertura-sub} confirma la huella territorial de este
sistema. El mapa evidencia cómo la cobertura del Tren Suburbano apenas alcanza la
alcaldía Cuauhtémoc y el extremo norte de Azcapotzalco, mientras que el resto del
territorio capitalino permanece fuera de su radio de influencia peatonal.

\begin{figure}[H]
    \centering
    \figinclude[width=0.85\textwidth]{Figures/Cap5/Cobertura/Mapas/NB02_fig02_SUB_cobertura.png}
    \caption[Mapa de cobertura territorial del Tren Suburbano en la CDMX]{Mapa
    del nivel de cobertura territorial del sistema Tren Suburbano por demarcación
    en la Ciudad de México, calculado a partir de \termino{buffers} de caminata
    peatonal de 800~m.}
    \label{fig:mapa-cobertura-sub}
    \fuentefigura{Elaboración propia con \texttt{VFTModel} utilizando visualización geoespacial.}
\end{figure}

\subsubsection{Tren Interurbano (INTERURBANO)}
\label{subsubsec:cobertura-interurbano}

El Tren Interurbano ``El Insurgente''\DIFdelbegin \DIFdel{es un sistema ferroviario de carácter
regional diseñado para conectar el Valle de Toluca con la Ciudad de México. Su
trazo dentro de la capital es incipiente y se concentra exclusivamente }\DIFdelend \DIFaddbegin \DIFadd{, de vocación regional Toluca-}\cdmx{}\DIFadd{,
concentra su presencia }\DIFaddend en el sector poniente \DIFdelbegin \DIFdel{, sirviendo como puerta de entrada a la metrópoli. Debido a su vocación
interurbana, su impacto en la cobertura espacial a nivel de alcaldías es altamente
localizado.
}%DIFDELCMD < 

%DIFDELCMD < %%%
\DIFdel{El Cuadro~\ref{tab:cobertura-interurbano-cdmx} presenta las alcaldías con mayor
nivel de cobertura para este sistema. Como reflejo de su ruta de ingreso, }\DIFdelend \DIFaddbegin \DIFadd{con cobertura efectiva }\DIFaddend únicamente \DIFaddbegin \DIFadd{en
}\DIFaddend Cuajimalpa de Morelos \DIFaddbegin \DIFadd{(2.23\,\%) }\DIFaddend y Álvaro Obregón \DIFdelbegin \DIFdel{presentan áreas de influencia dentro del
isócrono peatonal de 800~m, dejando al }\DIFdelend \DIFaddbegin \DIFadd{(0.44\,\%), mientras el }\DIFaddend resto
de la ciudad \DIFdelbegin \DIFdel{sin cobertura directa
por este medio}\DIFdelend \DIFaddbegin \DIFadd{carece de acceso directo al sistema
(Cuadro~\ref{tab:cobertura-interurbano-cdmx}). Para la visualización cartográfica
de este sistema en el contexto de la red metropolitana véase la
Figura~\ref{fig:anx-mapa-cobertura-total} del
Anexo~\ref{anx:cartografia-presentacion}}\DIFaddend .

\begin{table}[H]
    \centering
    \caption[Alcaldías con mayor cobertura de Tren Interurbano (CDMX)]{
    Cinco alcaldías de la Ciudad de México con mayor porcentaje de cobertura
    territorial para el sistema Tren Interurbano (INTERURBANO). Isócrono: 800~m.
    Corrida: 2026-04-18.}
    \label{tab:cobertura-interurbano-cdmx}
    \begin{tabular}{|l|r|r|r|}
        \hline
        \textbf{Demarcación} & \textbf{Área total (km²)} & \textbf{Área cubierta (km²)} & \textbf{Cobertura (\%)} \\
        \hline
        Cuajimalpa de Morelos &  71.10 & 1.59 & 2.23 \\
        Álvaro Obregón        &  95.82 & 0.42 & 0.44 \\
        Venustiano Carranza   &  33.84 & 0.00 & 0.00 \\
        Xochimilco            & 114.03 & 0.00 & 0.00 \\
        Tláhuac               &  85.78 & 0.00 & 0.00 \\
        \hline
    \end{tabular}
    \fuentefigura{Elaboración propia con \texttt{VFTModel} y datos de \texttt{Apimetro}.}
\end{table}

\DIFdelbegin \DIFdel{La Figura~\ref{fig:mapa-cobertura-interurbano} visualiza este alcance espacial
limitado. El mapa hace evidente la focalización de la infraestructura en el límite
occidental de la cuenca, evidenciando su papel como colector foráneo más que como
distribuidor urbano capilar.
}%DIFDELCMD < 

%DIFDELCMD < \begin{figure}[H]
%DIFDELCMD <     \centering
%DIFDELCMD <     \figinclude[width=0.85\textwidth]{Figures/Cap5/Cobertura/Mapas/NB02_fig02_INTERURBANO_cobertura.png}
%DIFDELCMD <     %%%
%DIFDELCMD < \caption[Mapa de cobertura territorial del Tren Interurbano en la CDMX]{%
{%DIFAUXCMD
\DIFdelFL{Mapa
    del nivel de cobertura territorial del sistema Tren Interurbano por
    demarcación en la Ciudad de México, calculado a partir de }%DIFDELCMD < \termino{buffers}
%DIFDELCMD <     %%%
\DIFdelFL{de caminata peatonal de 800~m.}}
    %DIFAUXCMD
%DIFDELCMD < \label{fig:mapa-cobertura-interurbano}
%DIFDELCMD <     \fuentefigura{Elaboración propia con \texttt{VFTModel} utilizando visualización geoespacial.}
%DIFDELCMD < \end{figure}
%DIFDELCMD < 

%DIFDELCMD < %%%
\DIFdelend \subsubsection{Pumabús (PB)}
\label{subsubsec:cobertura-pb}

El sistema Pumabús es un servicio de transporte interno operado por la Universidad
Nacional Autónoma de México (UNAM), de carácter público y gratuito, cuya operación
\DIFdelbegin \DIFdel{es exclusiva del }\DIFdelend \DIFaddbegin \DIFadd{se circunscribe al }\DIFaddend campus de Ciudad Universitaria. \DIFdelbegin \DIFdel{No obstante, debido a su relevancia
modal y estructural en el sur de la ciudad, se encuentra registrado formalmente
dentro del Portal de Datos Abiertos de la Ciudad de México, lo que permite integrar
su infraestructura de paradas y circuitos en el análisis macro de la conectividad
metropolitana.
}%DIFDELCMD < 

%DIFDELCMD < %%%
\DIFdel{El }\DIFdelend \DIFaddbegin \DIFadd{Si bien el modelo registra una
cobertura del 20.13\,\% en Coyoacán ---el porcentaje más alto entre los sistemas
de escala campus---, este valor refleja la concentración geométrica del isócrono
dentro del polígono universitario y no un alcance funcional sobre la trama urbana
circundante. El resto de las demarcaciones capitalinas registra coberturas menores
al 2.5\,\% o nulas, confirmando su carácter de sistema cautivo del recinto
universitario (}\DIFaddend Cuadro~\ref{tab:cobertura-pb-cdmx}\DIFdelbegin \DIFdel{identifica las cinco alcaldías de la Ciudad
de México con mayor cobertura territorial respecto a la red del Pumabús,
evidenciando su completa focalización en el núcleo universitario y su vecindad
inmediata}\DIFdelend \DIFaddbegin \DIFadd{). Para la visualización
cartográfica de este sistema en el contexto de la red metropolitana véase la
Figura~\ref{fig:anx-mapa-cobertura-total} del
Anexo~\ref{anx:cartografia-presentacion}}\DIFaddend .

\begin{table}[H]
    \centering
    \caption[Alcaldías con mayor cobertura de Pumabús (CDMX)]{
    Cinco alcaldías de la Ciudad de México con mayor porcentaje de cobertura
    territorial para el sistema Pumabús (PB). Isócrono: 800~m.
    Corrida: 2026-04-18.}
    \label{tab:cobertura-pb-cdmx}
    \begin{tabular}{|l|r|r|r|}
        \hline
        \textbf{Demarcación} & \textbf{Área total (km²)} & \textbf{Área cubierta (km²)} & \textbf{Cobertura (\%)} \\
        \hline
        Coyoacán            &  53.88 & 10.85 &  20.13 \\
        Álvaro Obregón      &  95.82 &  2.25 &   2.34 \\
        Tlalpan             & 314.25 &  0.25 &   0.08 \\
        Venustiano Carranza &  33.84 &  0.00 &   0.00 \\
        Xochimilco          & 114.03 &  0.00 &   0.00 \\
        \hline
    \end{tabular}
    \fuentefigura{Elaboración propia con \texttt{VFTModel} y datos de \texttt{Apimetro}.}
\end{table}


\DIFdelbegin \DIFdel{La configuración geoespacial de esta influencia se presenta en la Figura~\ref{fig:mapa-cobertura-pb}. El mapa hace patente la delimitación del
servicio al polígono escolar de Ciudad Universitaria, extendiendo su área de
influencia peatonal de 800~m únicamente hacia las franjas fronterizas de las tres
alcaldías circundantes}\DIFdelend %DIF >  =============================================================================
%DIF >   Vista cartográfica de la cobertura total de la red
%DIF >  =============================================================================
\DIFaddbegin \subsubsection{\DIFadd{Vista Cartográfica de la Cobertura Total}}
\label{subsubsec:cobertura-total-cartografia}

\DIFadd{El análisis pormenorizado por sistema presentado en las secciones anteriores
permite construir una lectura integral del territorio metropolitano en términos
de accesibilidad al transporte masivo. Las
Figuras~\ref{fig:mapa-cobertura-red-total} y~\ref{fig:mapa-cobertura-calor}
sintetizan esta lectura mediante dos representaciones cartográficas complementarias
producidas con el sistema }\textbf{\DIFadd{Transport-gis-zmvm-mjg}}\DIFadd{: la primera muestra el
alcance territorial acumulado de los isócronos de todos los sistemas; la segunda
expresa la densidad de esa cobertura como mapa de calor, revelando los núcleos de
máxima accesibilidad y los vacíos periféricos. Las versiones a página completa se
encuentran en el Anexo~\ref{anx:cartografia-presentacion}
(Figuras~\ref{fig:anx-mapa-cobertura-total}
y~\ref{fig:anx-mapa-cobertura-calor})}\DIFaddend .

\begin{figure}[H]
    \centering
    \DIFdelbeginFL %DIFDELCMD < \figinclude[width=0.85\textwidth]{Figures/Cap5/Cobertura/Mapas/NB02_fig02_PUMABUS_cobertura.png}
%DIFDELCMD <     \caption[Mapa de cobertura territorial del Pumabús en la CDMX]{%%%
\DIFdelFL{Mapa del
    nivel }\DIFdelendFL \DIFaddbeginFL \figinclude[width=0.90\textwidth]{Figures/Mapas/Mapa_Cobertura_Red_Total}
    \caption[Cobertura espacial total de la red multimodal — isócronas 800~m]{
    \DIFaddFL{Cobertura espacial integrada }\DIFaddendFL de \DIFdelbeginFL \DIFdelFL{cobertura territorial del sistema Pumabús por demarcación en la
    Ciudad }\DIFdelendFL \DIFaddbeginFL \DIFaddFL{los doce sistemas }\DIFaddendFL de \DIFdelbeginFL \DIFdelFL{México}\DIFdelendFL \DIFaddbeginFL \DIFaddFL{transporte masivo de la
    }\zmvm{}\DIFaddendFL , \DIFdelbeginFL \DIFdelFL{calculado }\DIFdelendFL \DIFaddbeginFL \DIFaddFL{calculada }\DIFaddendFL a partir de \DIFdelbeginFL %DIFDELCMD < \termino{buffers} %%%
\DIFdelendFL \DIFaddbeginFL \DIFaddFL{isócronos peatonales }\DIFaddendFL de \DIFdelbeginFL \DIFdelFL{caminata
    peatonal de }\DIFdelendFL 800~m. \DIFaddbeginFL \DIFaddFL{Las zonas sin
    cobertura se concentran en el sur y oriente de la cuenca metropolitana,
    evidenciando las demarcaciones con mayor exclusión territorial del sistema.}\DIFaddendFL }
    \DIFdelbeginFL %DIFDELCMD < \label{fig:mapa-cobertura-pb}
%DIFDELCMD <     \fuentefigura{Elaboración propia con \texttt{VFTModel} utilizando visualización geoespacial.}
%DIFDELCMD < %%%
\DIFdelendFL \DIFaddbeginFL \label{fig:mapa-cobertura-red-total}
    \fuentefigura{Elaboración propia con \texttt{Transport-gis-zmvm-mjg};
    datos: GTFS SEMOVI 2024 e INEGI 2023.}
\end{figure}

\begin{figure}[H]
    \centering
    \figinclude[width=0.90\textwidth]{Figures/Mapas/Mapa_Cobertura_Mapa_Calor}
    \caption[Densidad de cobertura de la red multimodal — mapa de calor]{
    \DIFaddFL{Densidad de la cobertura espacial de la red multimodal expresada como mapa
    de calor. Los núcleos de máxima accesibilidad se ubican en las alcaldías
    centrales (Cuauhtémoc, Benito Juárez, Iztacalco), mientras que las periferias
    sur y oriente presentan densidades cercanas a cero, confirmando la inequidad
    territorial diagnosticada.}}
    \label{fig:mapa-cobertura-calor}
    \fuentefigura{Elaboración propia con \texttt{Transport-gis-zmvm-mjg};
    datos: GTFS SEMOVI 2024 e INEGI 2023.}
\DIFaddendFL \end{figure}

%% ── 5-Resultados/5-2-1-Fase1/5-2-1-3-FuerzaCapilar.tex
\subsection*{5 2 1 3 FuerzaCapilar}
% !TeX root = ../../tesis.tex
% =============================================================================
%  5.2.1.3 FUERZA CAPILAR Y NODOS DE TRANSFERENCIA
% =============================================================================
\subsection{\bajoRevision{Fuerza Capilar y Nodos de Transferencia (\texorpdfstring{$FC$}{FC})}}
\label{subsec:fuerza-capilar}
\notaRevision{Los resultados de esta subsección están sujetos a revisión con el
comité tutor. Los valores presentados corresponden a la corrida del modelo del
18 de abril de 2026.}

La \textbf{Fuerza Capilar} ($FC$) cuantifica el nivel de convergencia de flujos
de transporte en cada nodo o agrupación de nodos del grafo. El indicador se
evalúa desde dos perspectivas complementarias:

\textbf{A. Fuerza Capilar Geográfica ($FC_H$) — Macro-Hubs:} agrupa
topológicamente estaciones físicamente colindantes (radio de 100~m) en un
subgrafo $H$ que consolida los flujos de sus estaciones constituyentes,
representando el fenómeno urbano del transbordo físico (CETRAMs):

\begin{equation}
    FC_{H} = \sum_{j \in H} \bigl(k_{\text{in}}(j) + k_{\text{out}}(j)\bigr)
    - \text{Aristas Internas}
    \label{eq:fc-geografica}
\end{equation}

El Cuadro~\ref{tab:fuerza-capilar-hubs} detalla el comportamiento numérico de los
cinco principales Macro-Hubs de transferencia identificados en la red con base en
esta formulación, mientras que la Figura~\ref{fig:barplot-macro-hubs} permite
contrastar las magnitudes relativas de la concentración de flujos en los diez
nodos más significativos de la metrópoli.

\begin{table}[H]
    \centering
    \caption[Top 5 Macro-Hubs con mayor Fuerza Capilar Geográfica ($FC_H$)]{
    Los cinco Macro-Hubs de transferencia (nodos agrupados topológicamente a
    menos de 100~m) con mayor Fuerza Capilar Geográfica en la red metropolitana.
    Corrida: 2026-04-18.}
    \label{tab:fuerza-capilar-hubs}
    \fontsize{8.5pt}{10.5pt}\selectfont
    \begin{tabularx}{\textwidth}{|c|X|c|c|c|c|}
        \hline
        \textbf{ID Hub} & \textbf{Principales Estaciones Agrupadas} & \textbf{Sistemas} & \textbf{Nodos} & \textbf{Conexiones} \textit{(In / Out)} & \textbf{$FC_H$} \\
        \hline
        2174 & Zona Hospitales Sur (San Fernando, Chimalcoyotl, Coapa, Hosp. Manuel Gea Gtz., INER, Urgencias Respiratorias) & CC, RTP & 21 & 66 / 64 & \textbf{130} \\
        \hline
        136 & Eje Huipulco -- Estadio Azteca (Acoxpa, Huipulco, Estadio Azteca, Renato Leduc, San Alejandro) & CC, RTP, TL & 61 & 61 / 50 & \textbf{111} \\
        \hline
        1597 & Corredor Tlalpan Sur (Calle La Rosa, Tezoquipa, Galeana, Jojutla) & CC, RTP & 10 & 49 / 51 & \textbf{100} \\
        \hline
        259 & Nodo Barranca del Muerto (Barranca del Muerto, Alpes, Periférico Sur, José María Velasco, Miguel Ocaranza) & CC, RTP & 21 & 50 / 42 & \textbf{92} \\
        \hline
        129 & CETRAM Chapultepec (Metro Chapultepec, Acapulco, Agustín Melgar, Sonora, Veracruz, Tampico) & CC, METRO, RTP, TROLE & 29 & 37 / 40 & \textbf{77} \\
        \hline
    \end{tabularx}
    \vspace{1ex}
    \fuentefigura{Elaboración propia con \texttt{VFTModel} y datos de \texttt{Apimetro}.}
\end{table}

\begin{figure}[H]
    \centering
    \figinclude[width=0.85\textwidth]{Figures/Cap5/FuerzaCapilar/Tabla/NB03_fig01_top10_macrohubs_barplot.png}
    \caption[Distribución de los principales Macro-Hubs por Fuerza Capilar]{
    Gráfico de barras comparativo de los diez principales Macro-Hubs de la red
    metropolitana, ordenados con base en su nivel acumulado de Fuerza Capilar
    Geográfica ($FC_H$).}
    \label{fig:barplot-macro-hubs}
    \fuentefigura{Elaboración propia con \texttt{VFTModel} a partir de los resultados analíticos del modelo.}
\end{figure}

Por último, la distribución espacial y la cobertura geográfica de estos nodos de
transferencia masiva dentro del entorno urbano de la \zmvm{} se presenta mapeada
de forma detallada en la Figura~\ref{fig:mapa-fuerza-capilar}.

\begin{figure}[H]
    \centering
    \figinclude[width=0.85\textwidth]{Figures/Cap5/FuerzaCapilar/Mapa/NB03_mapa01_macrohubs.png}
    \caption[Mapa de Macro-Hubs y Fuerza Capilar Geográfica]{
    Representación geoespacial de la \termino{Fuerza Capilar Geográfica}
    ($FC_H$). El mapa destaca la concentración de flujos de transporte en los
    principales nodos y zonas de transferencia física (Macro-Hubs) de la \zmvm.}
    \label{fig:mapa-fuerza-capilar}
    \fuentefigura{Elaboración propia con \texttt{VFTModel} utilizando visualización geoespacial.}
\end{figure}

% =============================================================================
%  5.2.1.4 FUERZA CAPILAR SISTÉMICA (MATEMÁTICA)
% =============================================================================
\subsection{\bajoRevision{Fuerza Capilar Sistémica (Matemática)}}
\label{subsec:fuerza-capilar-matematica}

\textbf{B. Fuerza Capilar Sistémica o Matemática ($FC$):} A diferencia de la
aproximación geográfica orientada a macro-hubs, esta perspectiva evalúa de manera
estricta la convergencia de flujos directamente sobre cada nodo topológico
individual. El análisis permite identificar los puntos exactos de la
infraestructura matemática de la red donde se concentra la mayor densidad de
aristas de entrada y salida, aislando singularidades críticas en la traza sin
aplicar tolerancias de proximidad espacial ni fusión de estaciones colindantes.

El Cuadro~\ref{tab:fuerza-capilar-matematica} detalla el comportamiento de los
cinco nodos individuales que registraron la mayor Fuerza Capilar pura dentro del
grafo de transporte metropolitano.

\begin{table}[H]
    \centering
    \caption[Top 5 Nodos con mayor Fuerza Capilar Sistémica (Matemática)]{
    Los cinco nodos individuales con mayor Fuerza Capilar Sistémica (Matemática)
    en la red metropolitana, evaluados de manera estricta sin agrupación
    geográfica. Corrida: 2026-04-18.}
    \label{tab:fuerza-capilar-matematica}
    \fontsize{9.5pt}{11.5pt}\selectfont
    \begin{tabular}{|c|l|c|c|c|}
        \hline
        \textbf{ID Nodo} & \textbf{Estación (Nodo Topológico)} & \textbf{Sistemas} & \textbf{Conexiones} \textit{(In / Out)} & \textbf{$FC$} \\
        \hline
        7270  & Luis Murillo                          & RTP & 21 / 20 & \textbf{41} \\
        \hline
        2817  & Calz. de Tlalpan -- Glorieta Huipulco & RTP & 19 / 15 & \textbf{34} \\
        \hline
        10211 & Sobre Calz. de Tlalpan y Luis Murillo & CC  & 16 / 17 & \textbf{33} \\
        \hline
        10717 & Tlalpan y Luis Murillo                & CC  & 18 / 15 & \textbf{33} \\
        \hline
        5540  & Estadio Azteca 2                      & RTP & 14 / 18 & \textbf{32} \\
        \hline
    \end{tabular}
    \vspace{1ex}
    \fuentefigura{Elaboración propia con \texttt{VFTModel} y datos de \texttt{Apimetro}.}
\end{table}

Las magnitudes y diferencias relativas entre los diez elementos individuales con
mayor concentración de conexiones se visualizan en el gráfico de barras de la
Figura~\ref{fig:barplot-nodos-matematicos}.

\begin{figure}[H]
    \centering
    \figinclude[width=0.85\textwidth]{Figures/Cap5/FuerzaCapilar/Tabla/NB03_fig02_top10_nodos_matematicos_barplot.png}
    \caption[Distribución de los principales nodos individuales por Fuerza Capilar]{
    Gráfico de barras comparativo de los diez nodos individuales con mayor
    Fuerza Capilar Sistémica (Matemática) en la red metropolitana.}
    \label{fig:barplot-nodos-matematicos}
    \fuentefigura{Elaboración propia con \texttt{VFTModel} a partir de los resultados analíticos del modelo.}
\end{figure}

La distribución geográfica e implicación espacial de estos puntos específicos de
alta carga conectiva dentro de la red urbana consolidada se presenta en la
Figura~\ref{fig:mapa-nodos-matematicos}.

\begin{figure}[H]
    \centering
    \figinclude[width=0.85\textwidth]{Figures/Cap5/FuerzaCapilar/Mapa/NB03_mapa02_nodos_matematicos.png}
    \caption[Mapa de Nodos Individuales y Fuerza Capilar Sistémica]{
    Representación geoespacial de la Fuerza Capilar Sistémica (Matemática) a
    nivel de nodo individual en la \zmvm{}.}
    \label{fig:mapa-nodos-matematicos}
    \fuentefigura{Elaboración propia con \texttt{VFTModel} utilizando visualización geoespacial.}
\end{figure}

\subsubsection{Análisis del Eje Periférico — Infraestructura Preexistente}
\label{subsubsec:periferico-fc}

Como diagnóstico previo a la propuesta del corredor anillar, se aislaron los nodos
y Macro-Hubs cuyos nombres contienen la denominación ``Periférico'', con el objetivo
de identificar los puntos de mayor tensión vehicular e intermodal sobre esta
vialidad fundamental. Este análisis se divide en dos enfoques: la concentración
geográfica (transbordos físicos en Macro-Hubs) y la concentración matemática
estricta (nodos topológicos individuales).

\paragraph{A. Fuerza Capilar Geográfica (Macro-Hubs) en el Periférico}

Al aplicar un radio de tolerancia de 100~m para agrupar las estaciones colindantes
sobre el eje, emergen los principales complejos de transbordo y alimentación. El
Cuadro~\ref{tab:periferico-hubs-geo} detalla los cinco Macro-Hubs de mayor Fuerza
Capilar Geográfica ($FC_H$), destacando la importancia de Barranca del Muerto y la
intersección con la zona de hospitales/Tlalpan como los anclajes más fuertes del
corredor.

\begin{table}[H]
    \centering
    \caption[Top 5 Macro-Hubs con mayor Fuerza Capilar sobre el Eje Periférico]{
    Cinco Macro-Hubs de transferencia con mayor Fuerza Capilar Geográfica
    ($FC_H$) ubicados sobre el Anillo Periférico. Corrida: 2026-04-18.}
    \label{tab:periferico-hubs-geo}
    \fontsize{8.5pt}{10.5pt}\selectfont
    \begin{tabularx}{\textwidth}{|c|X|c|c|c|c|}
        \hline
        \textbf{ID Hub} & \textbf{Principales Estaciones Agrupadas} & \textbf{Sistemas} & \textbf{Nodos} & \textbf{Conexiones} \textit{(In / Out)} & \textbf{$FC_H$} \\
        \hline
        259 & Nodo Barranca del Muerto (Barranca del Muerto, Alpes, Periférico Sur, Miguel Ocaranza) & CC, RTP & 21 & 50 / 42 & \textbf{92} \\
        \hline
        374 & Intersección Tlalpan--Periférico (Renato Leduc, Huipulco, Periférico, Av. Pedregal) & CC, RTP & 18 & 34 / 33 & \textbf{67} \\
        \hline
        453 & Nodo San Antonio (San Antonio, Periférico Sur, Puente Periférico Norte) & CC, RTP & 7 & 23 / 23 & \textbf{46} \\
        \hline
        402 & Eje Oriente Tepalcates (Constitución de Apatzingán, Canal de San Juan, Tepalcates) & CC, MB, RTP & 8 & 20 / 26 & \textbf{46} \\
        \hline
        291 & Intersección Xochimilco--Tepepan (16 de Septiembre, Kioto, Tec de Monterrey, Periférico Participación) & CC, RTP, TL & 14 & 23 / 22 & \textbf{45} \\
        \hline
    \end{tabularx}
    \vspace{1ex}
    \fuentefigura{Elaboración propia con \texttt{VFTModel} y datos de \texttt{Apimetro}.}
\end{table}

La Figura~\ref{fig:barplot-hubs-periferico} compara las magnitudes relativas de
estos agrupamientos geográficos, mientras que la Figura~\ref{fig:mapa-hubs-periferico}
espacializa su huella sobre la vialidad, revelando los polos naturales de atracción
de demanda que la futura infraestructura anillar deberá integrar.

\begin{figure}[H]
    \centering
    \figinclude[width=0.85\textwidth]{Figures/Cap5/FuerzaCapilar/Tabla/NB03_fig04_top10_macrohubs_periferico_barplot.png}
    \caption[Fuerza Capilar de Macro-Hubs en Anillo Periférico]{Gráfico de barras
    de los diez principales Macro-Hubs geográficos ubicados en el Anillo Periférico,
    ordenados por su nivel de Fuerza Capilar.}
    \label{fig:barplot-hubs-periferico}
    \fuentefigura{Elaboración propia con \texttt{VFTModel}.}
\end{figure}

\begin{figure}[H]
    \centering
    \figinclude[width=0.85\textwidth]{Figures/Cap5/FuerzaCapilar/Mapa/NB03_mapa03_macrohubs_periferico.png}
    \caption[Distribución espacial de Macro-Hubs en el Periférico]{Mapa de la
    Fuerza Capilar Geográfica ($FC_H$) en los Macro-Hubs del Anillo Periférico.}
    \label{fig:mapa-hubs-periferico}
    \fuentefigura{Elaboración propia con \texttt{VFTModel} utilizando visualización geoespacial.}
\end{figure}

\paragraph{B. Fuerza Capilar Sistémica (Nodos Matemáticos) en el Periférico}

Al desagregar los Hubs y analizar la Fuerza Capilar Sistémica ($FC$) estricta a
nivel de nodo individual, es posible identificar los paraderos específicos que
soportan la mayor carga de conexiones en la infraestructura actual. El
Cuadro~\ref{tab:periferico-nodos-math} presenta los cinco paraderos matemáticos
de mayor convergencia de flujo.

\begin{table}[H]
    \centering
    \caption[Nodos con mayor Fuerza Capilar Sistémica sobre el Eje Periférico]{
    Cinco nodos topológicos individuales con mayor Fuerza Capilar Matemática
    ($FC$) ubicados sobre el Anillo Periférico. Corrida: 2026-04-18.}
    \label{tab:periferico-nodos-math}
    \fontsize{9.5pt}{11.5pt}\selectfont
    \begin{tabular}{|c|l|c|r|r|c|}
        \hline
        \textbf{ID Nodo} & \textbf{Estación (Nodo Topológico)} & \textbf{Sistema} & \textbf{Entr.} & \textbf{Sal.} & \textbf{$FC$} \\
        \hline
        707  & Anillo Periférico y Blvd. Adolfo Ruiz Cortines & RTP & 14 & 16 & \textbf{30} \\
        \hline
        706  & Anillo Periférico y Blvd. Adolfo Ruiz Cortines & RTP & 13 & 13 & \textbf{26} \\
        \hline
        8538 & Periférico -- Barranca del Muerto              &  CC & 12 &  9 & \textbf{21} \\
        \hline
        602  & Anillo Periférico Canal de Garay -- Bilbao     & RTP & 10 & 10 & \textbf{20} \\
        \hline
        8510 & Periférico                                     & RTP &  9 & 11 & \textbf{20} \\
        \hline
    \end{tabular}
    \vspace{1ex}
    \fuentefigura{Elaboración propia con \texttt{VFTModel} y datos de \texttt{Apimetro}.}
\end{table}

La jerarquía analítica de estos puntos singulares se constata en la
Figura~\ref{fig:barplot-nodos-periferico}, seguida de la
Figura~\ref{fig:mapa-nodos-periferico}, que detalla la distribución geográfica a
escala de paradero individual, útil para el micro-posicionamiento de futuras
estaciones anillares.

\begin{figure}[H]
    \centering
    \figinclude[width=0.85\textwidth]{Figures/Cap5/FuerzaCapilar/Tabla/NB03_fig05_top10_nodos_matematicos_periferico_barplot.png}
    \caption[Fuerza Capilar de nodos individuales en Anillo Periférico]{
    Gráfico de barras comparativo de los diez principales nodos matemáticos
    individuales ubicados en el Anillo Periférico.}
    \label{fig:barplot-nodos-periferico}
    \fuentefigura{Elaboración propia con \texttt{VFTModel}.}
\end{figure}

\begin{figure}[H]
    \centering
    \figinclude[width=0.85\textwidth]{Figures/Cap5/FuerzaCapilar/Mapa/NB03_mapa04_nodos_matematicos_periferico.png}
    \caption[Distribución espacial de nodos matemáticos en el Periférico]{Mapa
    de la Fuerza Capilar Sistémica (Matemática) a nivel de nodo individual sobre
    el corredor del Anillo Periférico.}
    \label{fig:mapa-nodos-periferico}
    \fuentefigura{Elaboración propia con \texttt{VFTModel} utilizando visualización geoespacial.}
\end{figure}

\vspace{2ex}
\begin{quote}
    \small\textit{Nota aclaratoria:} los datos presentados en esta sección
    reflejan exclusivamente la infraestructura de transporte \textbf{preexistente}
    sobre el Anillo Periférico. No se ha trazado ni modelado aún la geometría
    del corredor anillar propuesto; la sección permanece en plano especulativo
    para identificar los puntos naturales de anclaje de la futura línea a nivel
    capilar.
\end{quote}
\DIFaddbegin 

\DIFadd{La Figura~\ref{fig:mapa-fc-presentacion} consolida, mediante cartografía de
presentación, la distribución geoespacial de los principales nodos }\termino{hub}
\DIFadd{de la red y su intensidad capilar. La versión a página completa se encuentra en
la Figura~\ref{fig:anx-mapa-fuerza-capilar} del
Anexo~\ref{anx:cartografia-presentacion}.
}

\begin{figure}[H]
    \centering
    \figinclude[width=0.90\textwidth]{Figures/Mapas/Mapa_Fuerza_Capilar}
    \caption[Cartografía de presentación — Fuerza Capilar y nodos \textit{hub} principales]{
    \DIFaddFL{Representación cartográfica de los veinte nodos con mayor Fuerza Capilar en
    la red de transporte de la }\zmvm{}\DIFaddFL{. Los nodos señalados corresponden a los
    puntos de mayor influencia redistributiva de flujos intermodales, con especial
    concentración en el corredor sur de la cuenca metropolitana, sobre el eje
    que recorre el Anillo Periférico.}}
    \label{fig:mapa-fc-presentacion}
    \fuentefigura{Elaboración propia con \texttt{Transport-gis-zmvm-mjg};
    datos: GTFS SEMOVI 2024 e INEGI 2023.}
\end{figure}
 \DIFaddend

%% ── 5-Resultados/5-2-1-Fase1/5-2-1-4-FactorDesviacion.tex
\subsection*{5 2 1 4 FactorDesviacion}
% !TeX root = ../../tesis.tex
% =============================================================================
%  5.2.1.5 FACTOR DE DESVIACIÓN
% =============================================================================
\subsection{\bajoRevision{Factor de Desviación (\texorpdfstring{$DF$}{DF})}}
\label{subsec:detour-factor}
\notaRevision{Los resultados de esta subsección están sujetos a revisión con el
comité tutor. Los valores presentados corresponden a una muestra de 100~rutas
aleatorias ejecutadas el 27 de abril de 2026.}

El \textbf{Factor de Desviación} ($DF$) mide la eficiencia topológica de la
red comparando la distancia real recorrida a través de la infraestructura
contra la distancia teórica en línea recta (Haversine) entre el mismo par
origen-destino:

\begin{equation}
    DF = \frac{D_{\text{red}} + D_{\text{caminata}}}{D_{\text{haversine}}(O, D)}
    \label{eq:detour-factor-res}
\end{equation}

\begin{itemize}
    \item $D_{\text{red}}$: distancia acumulada sobre el trazo real de la ruta
    óptima en tiempo.
    \item $D_{\text{caminata}}$: suma de la primera y última milla peatonal.
    \item $D_{\text{haversine}}(O,D)$: distancia en línea recta entre origen y
    destino calculada con la fórmula de Haversine.
\end{itemize}

Un valor de $DF = 1.0$ indica una ruta perfectamente directa; valores superiores
a $1.5$ revelan ineficiencias topológicas significativas.

% -----------------------------------------------------------------------------
\subsubsection{Estadísticas Globales de la Red}
\label{subsubsec:df-estadisticas}

Sobre una muestra aleatoria de \textbf{100 pares origen-destino} (semilla
reproducible $= 42$), el análisis arrojó los siguientes estadísticos globales:

\begin{table}[H]
    \centering
    \caption[Estadísticos globales del Factor de Desviación — muestra de 100 rutas]{
    Estadísticos descriptivos del $DF$ sobre 100 pares O-D de la red de la \zmvm.
    Corrida: 2026-04-27.}
    \label{tab:df-estadisticas}
    \begin{tabular}{|l|r|}
        \hline
        \textbf{Estadístico} & \textbf{Valor} \\
        \hline
        $n$ (rutas calculadas)     & 100   \\
        Media ($\bar{DF}$)         & 1.545 \\
        Desviación estándar        & 0.517 \\
        Mínimo                     & 0.97  \\
        Percentil 25               & 1.27  \\
        Mediana                    & 1.40  \\
        Percentil 75               & 1.67  \\
        Máximo                     & 4.34  \\
        \hline
    \end{tabular}
    \fuentefigura{Elaboración propia con \texttt{VFTModel} (corrida 2026-04-27).}
\end{table}

La Figura~\ref{fig:df-caso-inicial} muestra la distribución geoespacial de los
100 pares O-D y las rutas calculadas sobre el territorio de la \zmvm{}.

\begin{figure}[H]
    \centering
    \figinclude[width=0.85\textwidth]{Figures/Cap5/FactorDesviacion/Mapa/NB04_mapa01_caso_inicial.png}
    \caption[Distribución de la muestra O-D para el Factor de Desviación]{
    Distribución de los 100 pares origen-destino de la muestra aleatoria y sus
    rutas óptimas en la \zmvm. Corrida: 2026-04-27.}
    \label{fig:df-caso-inicial}
    \fuentefigura{Elaboración propia con \texttt{VFTModel} (corrida 2026-04-27).}
\end{figure}

% -----------------------------------------------------------------------------
\subsubsection{Casos de Estudio Seleccionados}
\label{subsubsec:df-casos}

Se analizaron doce rutas que evalúan la conectividad desde las periferias sur y
oriente hacia distintos destinos de la metrópoli. Las
Figuras~\ref{fig:df-casos-1-4}--\ref{fig:df-casos-9-12} presentan los trazados
cartográficos de cada caso; el Cuadro~\ref{tab:df-casos} consolida los valores
numéricos al final de esta sección.

% --- Casos 1–4 ---------------------------------------------------------------
\begin{figure}[H]
    \centering
    \begin{subfigure}[t]{0.48\textwidth}
        \figinclude[width=\textwidth]{Figures/Cap5/FactorDesviacion/Mapa/NB04_mapa04_caso01.png}
        \caption{Caso~1: Milpa Alta → Autódromo. $DF = 1.29$.}
        \label{fig:df-c01}
    \end{subfigure}
    \hfill
    \begin{subfigure}[t]{0.48\textwidth}
        \figinclude[width=\textwidth]{Figures/Cap5/FactorDesviacion/Mapa/NB04_mapa05_caso02.png}
        \caption{Caso~2: Milpa Alta → Reforma (CBD). $DF = 1.37$.}
        \label{fig:df-c02}
    \end{subfigure}

    \vspace{0.5em}
    \begin{subfigure}[t]{0.48\textwidth}
        \figinclude[width=\textwidth]{Figures/Cap5/FactorDesviacion/Mapa/NB04_mapa06_caso03.png}
        \caption{Caso~3: Milpa Alta → Santa Fe (CBD). $DF = 1.27$.}
        \label{fig:df-c03}
    \end{subfigure}
    \hfill
    \begin{subfigure}[t]{0.48\textwidth}
        \figinclude[width=\textwidth]{Figures/Cap5/FactorDesviacion/Mapa/NB04_mapa07_caso04.png}
        \caption{Caso~4: Milpa Alta → FES Acatlán. $DF = 1.35$.}
        \label{fig:df-c04}
    \end{subfigure}
    \caption[Casos de estudio del Factor de Desviación: 1–4]{Rutas de los casos 1 al 4.
    Origen común: Milpa Alta hacia distintos destinos del norte y poniente.}
    \label{fig:df-casos-1-4}
    \fuentefigura{Elaboración propia con \texttt{VFTModel} (corrida 2026-04-27).}
\end{figure}

% --- Casos 5–8 ---------------------------------------------------------------
\begin{figure}[H]
    \centering
    \begin{subfigure}[t]{0.48\textwidth}
        \figinclude[width=\textwidth]{Figures/Cap5/FactorDesviacion/Mapa/NB04_mapa08_caso05.png}
        \caption{Caso~5: Milpa Alta → C.U. $DF = 1.32$.}
        \label{fig:df-c05}
    \end{subfigure}
    \hfill
    \begin{subfigure}[t]{0.48\textwidth}
        \figinclude[width=\textwidth]{Figures/Cap5/FactorDesviacion/Mapa/NB04_mapa09_caso06.png}
        \caption{Caso~6: Milpa Alta → Coyoacán. $DF = 1.34$.}
        \label{fig:df-c06}
    \end{subfigure}

    \vspace{0.5em}
    \begin{subfigure}[t]{0.48\textwidth}
        \figinclude[width=\textwidth]{Figures/Cap5/FactorDesviacion/Mapa/NB04_mapa10_caso07.png}
        \caption{Caso~7: Santa Catarina → C.U. $DF = 1.82$.}
        \label{fig:df-c07}
    \end{subfigure}
    \hfill
    \begin{subfigure}[t]{0.48\textwidth}
        \figinclude[width=\textwidth]{Figures/Cap5/FactorDesviacion/Mapa/NB04_mapa11_caso08.png}
        \caption{Caso~8: Santa Catarina → FES Acatlán. $DF = 1.42$.}
        \label{fig:df-c08}
    \end{subfigure}
    \caption[Casos de estudio del Factor de Desviación: 5–8]{Rutas de los casos 5 al 8.
    Destinos en el sur de la ciudad y primeros casos desde Santa Catarina.}
    \label{fig:df-casos-5-8}
    \fuentefigura{Elaboración propia con \texttt{VFTModel} (corrida 2026-04-27).}
\end{figure}

% --- Casos 9–12 --------------------------------------------------------------
\begin{figure}[H]
    \centering
    \begin{subfigure}[t]{0.48\textwidth}
        \figinclude[width=\textwidth]{Figures/Cap5/FactorDesviacion/Mapa/NB04_mapa12_caso09.png}
        \caption{Caso~9: Milpa Alta → Santa Catarina. $DF = 2.52$.}
        \label{fig:df-c09}
    \end{subfigure}
    \hfill
    \begin{subfigure}[t]{0.48\textwidth}
        \figinclude[width=\textwidth]{Figures/Cap5/FactorDesviacion/Mapa/NB04_mapa13_caso10.png}
        \caption{Caso~10: Milpa Alta → Ajusco. $DF = 1.42$.}
        \label{fig:df-c10}
    \end{subfigure}

    \vspace{0.5em}
    \begin{subfigure}[t]{0.48\textwidth}
        \figinclude[width=\textwidth]{Figures/Cap5/FactorDesviacion/Mapa/NB04_mapa14_caso11.png}
        \caption{Caso~11: Ajusco → Preparatoria 1. $DF = 1.43$.}
        \label{fig:df-c11}
    \end{subfigure}
    \hfill
    \begin{subfigure}[t]{0.48\textwidth}
        \figinclude[width=\textwidth]{Figures/Cap5/FactorDesviacion/Mapa/NB04_mapa15_caso12.png}
        \caption{Caso~12: P. Los Olivos → Centro Tláhuac. $DF = 1.06$.}
        \label{fig:df-c12}
    \end{subfigure}
    \caption[Casos de estudio del Factor de Desviación: 9–12]{Rutas de los casos 9 al 12.
    El caso~9 (sur-oriente transversal) registra el mayor $DF$ del conjunto.}
    \label{fig:df-casos-9-12}
    \fuentefigura{Elaboración propia con \texttt{VFTModel} (corrida 2026-04-27).}
\end{figure}

% --- Extremos globales de la muestra aleatoria -------------------------------
Para situar los casos de estudio en el contexto de la muestra completa, las
Figuras~\ref{fig:df-mejor-caso} y~\ref{fig:df-peor-caso} ilustran los extremos
globales obtenidos en los 100 pares O-D aleatorios.

\begin{figure}[H]
    \centering
    \begin{subfigure}[t]{0.48\textwidth}
        \figinclude[width=\textwidth]{Figures/Cap5/FactorDesviacion/Mapa/NB04_mapa02_ruta_mas_eficiente.png}
        \caption{Mejor caso: Plutarco Elías Calles → Fuente de Loreto. $DF = 0.97$.}
        \label{fig:df-mejor-caso}
    \end{subfigure}
    \hfill
    \begin{subfigure}[t]{0.48\textwidth}
        \figinclude[width=\textwidth]{Figures/Cap5/FactorDesviacion/Mapa/NB04_mapa03_ruta_mas_ineficiente.png}
        \caption{Peor caso: Carlos A. Vidal → 20 de Noviembre. $DF = 4.34$.}
        \label{fig:df-peor-caso}
    \end{subfigure}
    \caption[Extremos del Factor de Desviación en la muestra aleatoria]{
    Rutas con menor y mayor Factor de Desviación en la muestra aleatoria de 100 pares O-D.}
    \label{fig:df-extremos}
    \fuentefigura{Elaboración propia con \texttt{VFTModel} (corrida 2026-04-27).}
\end{figure}

% --- Tabla resumen de los 12 casos -------------------------------------------
\begin{table}[H]
    \centering
    \caption[Factor de Desviación — resumen de los doce casos de estudio]{
    Resultados del Factor de Desviación para los doce casos estratégicos de la
    \zmvm{} ilustrados en las figuras anteriores. Corrida: 2026-04-27.}
    \label{tab:df-casos}
    \small
    \begin{tabularx}{\textwidth}{|l|X|X|r|r|r|}
        \hline
        \textbf{\#} & \textbf{Origen} & \textbf{Destino} &
        \textbf{$D_{\text{red}}$ (km)} & \textbf{$D_{\text{Hav.}}$ (km)} & \textbf{$DF$} \\
        \hline
         1 & Milpa Alta        & Autódromo / Centro-Oriente   & 29.18 & 22.59 & 1.29 \\
         2 & Milpa Alta        & Paseo de la Reforma (CBD)    & 37.73 & 27.50 & 1.37 \\
         3 & Milpa Alta        & Santa Fe (CBD)               & 36.42 & 28.62 & 1.27 \\
         4 & Milpa Alta        & FES Acatlán                  & 48.34 & 35.73 & 1.35 \\
         5 & Milpa Alta        & Ciudad Universitaria         & 26.25 & 19.96 & 1.32 \\
         6 & Milpa Alta        & Coyoacán                     & 25.24 & 18.86 & 1.34 \\
         7 & Santa Catarina    & Ciudad Universitaria         & 42.06 & 23.10 & 1.82 \\
         8 & Santa Catarina    & FES Acatlán                  & 45.64 & 32.16 & 1.42 \\
         9 & Milpa Alta        & Santa Catarina (Sur-Oriente) & 37.22 & 14.79 & 2.52 \\
        10 & Milpa Alta        & Ajusco                       & 29.19 & 20.56 & 1.42 \\
        11 & Ajusco            & Preparatoria 1               & 15.91 & 11.16 & 1.43 \\
        12 & Parque Los Olivos & Centro de Tláhuac            &  2.20 &  2.08 & 1.06 \\
        \hline
        \multicolumn{5}{|l|}{\textbf{Mejor caso global (muestra aleatoria):}
        Plutarco Elias Calles → Fuente de Loreto} & \textbf{0.97} \\
        \multicolumn{5}{|l|}{\textbf{Peor caso global (muestra aleatoria):}
        Carlos A. Vidal → 20 de Noviembre} & \textbf{4.34} \\
        \hline
    \end{tabularx}
    \fuentefigura{Elaboración propia con \texttt{VFTModel} (corrida 2026-04-27).}
\end{table}

El análisis revela que los viajes periféricos sur-oriente (Caso~9: $DF = 2.52$,
Figura~\ref{fig:df-c09}) registran la mayor ineficiencia entre los casos de
estudio, obligando al usuario a recorrer 37.2~km de red para cubrir una distancia
lineal de 14.8~km. Esta situación responde a la ausencia de conectividad
transversal en el arco sur-oriente de la ciudad, que impone al tránsito
perimetral el ingreso al centro metropolitano como nodo obligado, confirmando la
hipótesis sobre la necesidad de un corredor anillar en el Periférico.
\DIFaddbegin 

\DIFadd{La Figura~\ref{fig:mapa-df-presentacion} sintetiza estos hallazgos mediante
cartografía de presentación del Factor de Desviación sobre el conjunto de nodos
de la red. El gradiente de color verde-a-rojo permite identificar directamente las
zonas donde la infraestructura existente impone los mayores rodeos, reforzando de
forma visual el argumento sobre la necesidad del corredor anillar en el arco
sur-oriente. La versión a página completa se encuentra en la
Figura~\ref{fig:anx-mapa-factor-desviacion} del
Anexo~\ref{anx:cartografia-presentacion}.
}

\begin{figure}[H]
    \centering
    \figinclude[width=0.90\textwidth]{Figures/Mapas/Mapa_DF}
    \caption[Cartografía de presentación — Factor de Desviación por nodo]{
    \DIFaddFL{Factor de Desviación ($DF$) por nodo en la red de transporte de la }\zmvm{}\DIFaddFL{.
    El gradiente verde-a-rojo indica la relación entre la distancia real de viaje
    y la distancia euclidiana: los nodos en rojo señalan las zonas donde las
    rutas disponibles imponen desviaciones significativas, penalizando con mayor
    intensidad a la periferia sur y oriente de la cuenca metropolitana.}}
    \label{fig:mapa-df-presentacion}
    \fuentefigura{Elaboración propia con \texttt{Transport-gis-zmvm-mjg};
    datos: GTFS SEMOVI 2024 e INEGI 2023.}
\end{figure}
 \DIFaddend

%% ── Anexos/E-Cartografia-Presentacion.tex
\subsection*{E Cartografia Presentacion}
% !TeX root = ../tesis.tex
%%%%%%%%%%%%%%%%%%%%%%%%%%%%%%%%%%%%%%%%%%%%%%%%%%%%%%%%%%%%%%%%%%%
% ANEXO E — Cartografía de Presentación: Transport-gis-zmvm-mjg
%%%%%%%%%%%%%%%%%%%%%%%%%%%%%%%%%%%%%%%%%%%%%%%%%%%%%%%%%%%%%%%%%%%

Este anexo reúne las versiones cartográficas de alta resolución producidas por el sistema \textbf{Transport-gis-zmvm-mjg},
generadas con QGIS~3.x a partir de las plantillas institucionales de la \textsc{unam}. Los mapas utilizan la proyección
\textsc{epsg}:32614 (UTM zona~14N) y tienen como fuente base los \textit{feeds} \gtfs{} de la \textsc{semovi} (2024) y la cartografía
del \textsc{inegi} (2023). Versiones reducidas de estos mapas se encuentran en los capítulos~\ref{ch:Capas-codigo}
y~\ref{ch:resultados-vft} para su discusión contextual.

% =============================================================================
\section{Red Esquemática Multimodal}
\label{anx:mapa-red-esquematica}
% =============================================================================

\begin{figure}[p]
  \centering
  \figinclude[width=\textwidth]{Figures/Mapas/Mapa_Red_Esquematica}
  \caption[Red esquemática multimodal de la ZMVM (alta resolución)]{Red esquemática multimodal de la \zmvm{}.
    Muestra los doce sistemas de transporte reconocidos por el modelo VFT sobre los límites político-administrativos
    de la \cdmx{} y el Estado de México.\DIFdelbeginFL %DIFDELCMD < \fuentefigura{Elaboración propia con \textbf{Transport-gis-zmvm-mjg};
%DIFDELCMD <     datos: \textsc{gtfs} \textsc{semovi} 2024 e \textsc{inegi} 2023.}%%%
\DIFdelendFL }
  \label{fig:anx-mapa-red}
  \DIFaddbeginFL \fuentefigura{Elaboración propia con \textbf{Transport-gis-zmvm-mjg};
    datos: GTFS SEMOVI 2024 e INEGI 2023.}
\DIFaddendFL \end{figure}

% =============================================================================
\section{Cobertura Espacial — Isócronas 800~m}
\label{anx:mapa-cobertura-total}
% =============================================================================

\begin{figure}[p]
  \centering
  \figinclude[width=\textwidth]{Figures/Mapas/Mapa_Cobertura_Red_Total}
  \caption[Cobertura de la red por isócronas de 800~m (alta resolución)]{Cobertura espacial de la red multimodal de transporte
    de la \zmvm{}, calculada como el área alcanzable a 800~m de caminata desde cualquier parada o estación
    (indicador $C$ del modelo VFT). Las alcaldías del sur presentan las coberturas más bajas.\DIFdelbeginFL %DIFDELCMD < \fuentefigura{Elaboración propia con \textbf{Transport-gis-zmvm-mjg};
%DIFDELCMD <     datos: \textsc{gtfs} \textsc{semovi} 2024 e \textsc{inegi} 2023.}%%%
\DIFdelendFL }
  \label{fig:anx-mapa-cobertura-total}
  \DIFaddbeginFL \fuentefigura{Elaboración propia con \textbf{Transport-gis-zmvm-mjg};
    datos: GTFS SEMOVI 2024 e INEGI 2023.}
\DIFaddendFL \end{figure}

% =============================================================================
\section{Cobertura Espacial — Mapa de Calor}
\label{anx:mapa-cobertura-calor}
% =============================================================================

\begin{figure}[p]
  \centering
  \figinclude[width=\textwidth]{Figures/Mapas/Mapa_Cobertura_Mapa_Calor}
  \caption[Densidad de cobertura — mapa de calor (alta resolución)]{Densidad de la cobertura espacial de la red
    representada como mapa de calor. La intensidad del color refleja la concentración de isócronas superpuestas,
    revelando los núcleos de mayor accesibilidad modal en la región metropolitana.\DIFdelbeginFL %DIFDELCMD < \fuentefigura{Elaboración propia con \textbf{Transport-gis-zmvm-mjg};
%DIFDELCMD <     datos: \textsc{gtfs} \textsc{semovi} 2024 e \textsc{inegi} 2023.}%%%
\DIFdelendFL }
  \label{fig:anx-mapa-cobertura-calor}
  \DIFaddbeginFL \fuentefigura{Elaboración propia con \textbf{Transport-gis-zmvm-mjg};
    datos: GTFS SEMOVI 2024 e INEGI 2023.}
\DIFaddendFL \end{figure}

% =============================================================================
\section{Fuerza Capilar — Nodos \textit{Hub} Principales}
\label{anx:mapa-fuerza-capilar}
% =============================================================================

\begin{figure}[p]
  \centering
  \figinclude[width=\textwidth]{Figures/Mapas/Mapa_Fuerza_Capilar}
  \caption[Fuerza capilar: los 20 nodos \textit{hub} de mayor centralidad (alta resolución)]{Fuerza capilar de los nodos de la red
    de transporte de la \zmvm{}. Se destacan los 20 nodos con mayor valor de centralidad de intermediación ponderada,
    que corresponden a los puntos de mayor influencia redistributiva de flujos en la red multimodal.\DIFdelbeginFL %DIFDELCMD < \fuentefigura{Elaboración propia con \textbf{Transport-gis-zmvm-mjg};
%DIFDELCMD <     datos: \textsc{gtfs} \textsc{semovi} 2024 e \textsc{inegi} 2023.}%%%
\DIFdelendFL }
  \label{fig:anx-mapa-fuerza-capilar}
  \DIFaddbeginFL \fuentefigura{Elaboración propia con \textbf{Transport-gis-zmvm-mjg};
    datos: GTFS SEMOVI 2024 e INEGI 2023.}
\DIFaddendFL \end{figure}

% =============================================================================
\section{Factor de Desviación por Nodo}
\label{anx:mapa-factor-desviacion}
% =============================================================================

\begin{figure}[p]
  \centering
  \figinclude[width=\textwidth]{Figures/Mapas/Mapa_DF}
  \caption[Factor de desviación por nodo (alta resolución)]{Factor de desviación ($\text{FD}$) de los nodos de la red
    de transporte de la \zmvm{}. El gradiente de color (verde→rojo) indica la relación entre la distancia real de viaje
    y la distancia euclidiana óptima: valores altos (rojo) señalan nodos donde las rutas disponibles desvían significativamente
    del trayecto directo, lo que penaliza la eficiencia de la red en esas zonas.\DIFdelbeginFL %DIFDELCMD < \fuentefigura{Elaboración propia con \textbf{Transport-gis-zmvm-mjg};
%DIFDELCMD <     datos: \textsc{gtfs} \textsc{semovi} 2024 e \textsc{inegi} 2023.}%%%
\DIFdelendFL }
  \label{fig:anx-mapa-factor-desviacion}
  \DIFaddbeginFL \fuentefigura{Elaboración propia con \textbf{Transport-gis-zmvm-mjg};
    datos: GTFS SEMOVI 2024 e INEGI 2023.}
\DIFaddendFL \end{figure}
